% concatenated from main.tex
 \documentclass[12pt,reqno]{amsart}
  
%\usepackage[active]{srcltx} %for YandY Tex

%%%% DELETE THESE LINES FOR FINAL VERSION
\newcommand{\ver}{\it{version 20}}%<---PUT VERSION NUMBER HERE
\newcommand{\cs}{$\clubsuit$} % THIS IS TO MARK COMMENTS, CHANGES
%\usepackage{showkeys}  % THIS SHOWS LABELS
%%%%  
 \usepackage{comment}
\headheight=6.15pt \textheight=8.75in \textwidth=6.5in
\oddsidemargin=0in \evensidemargin=0in \topmargin=0in
    
\usepackage{tikz}
\usepackage{mathrsfs}
\usepackage{latexsym}
\usepackage{graphicx}
\usepackage{amssymb}
 \graphicspath{{Figures/}}
 \usepackage[show]{ed}


\renewcommand{\Re}{{\operatorname{Re}\,}}
\renewcommand{\Im}{{\operatorname{Im}\,}}
\newcommand{\Op}{\operatorname{Op}}
\newcommand{\sgn}{\operatorname{sgn}}
\newcommand{\lla}{\left\langle}
\newcommand{\rra}{\right\rangle}
\newcommand{\RP}{\R\PP}
\newcommand{\bbb}{|\!|\!|}
\newcommand{\s}{{\mathbf s}}
\renewcommand{\epsilon}{\varepsilon}
\newcommand{\dist}{{\operatorname{dist}}}
\newcommand{\var}{{\operatorname{Var}}}
\newcommand{\Var}{{\bf Var}}
\newcommand{\poly}{{\operatorname{Poly}}}
\newcommand{\I}{{\mathbf I}}
\newcommand{\todd}{{\operatorname{Todd}}}
\newcommand{\T}{{\mathbf T}^m}
\newcommand{\boxb}{\raisebox{-.8pt}{$\Box$}_b}
\newcommand{\sm}{\smallsetminus}
\newcommand{\szego}{Szeg\H{o} }
\newcommand{\nhat}{\raisebox{2pt}{$\wh{\ }$}}
\newcommand{\Si}{\Sigma}
\newcommand{\inv}{^{-1}}
\newcommand{\kahler}{K\"ahler }
\newcommand{\sqrtn}{\sqrt{N}}

\newcommand{\Vt}{\Vert}
\newcommand{\wh}{\widehat}
\newcommand{\PP}{{\mathbb P}}
\newcommand{\N}{{\mathbb N}}
\newcommand{\R}{{\mathbb R}}
\newcommand{\C}{{\mathbb C}}
\newcommand{\Q}{{\mathbb Q}}
\newcommand{\Z}{{\mathbb Z}}
\newcommand{\B}{{\mathbb B}}
\newcommand{\Ss}{{\mathbb S}}
\newcommand{\Hh}{{\mathbb H}}
\newcommand{\lan}{\left\langle}
\newcommand{\ran}{\right\rangle}
\newcommand{\mc}[1]{\mathcal{#1}}
\newcommand{\e}{\epsilon}
\newcommand{\imply}{\Rightarrow}
\newcommand{\re}{\mathbb{R}}
\newcommand{\Cc}{C_c^\infty}
\newcommand{\Id}{\operatorname{Id}}

\newcommand{\h}{\hbar}
\newcommand{\CP}{\C\PP}
\renewcommand{\d}{\partial}
\newcommand{\dbar}{\bar\partial}
\newcommand{\ddbar}{\partial\dbar}
\newcommand{\U}{{\rm U}}
\newcommand{\D}{{\mathbf D}}
\newcommand{\K}{{\mathbf K}}
\newcommand{\E}{{\mathbf E}}
\newcommand{\Heis}{{\mathbf H}}
\newcommand{\zb}{{\mathbf Z}}
\newcommand{\Log}{{\operatorname{Log\,}}}
\newcommand{\half}{{\textstyle \frac 12}}
\newcommand{\vol}{{\operatorname{Vol}}}
\newcommand{\diag}{{\operatorname{diag}}}
\newcommand{\codim}{{\operatorname{codim\,}}}
\newcommand{\SU}{{\operatorname{SU}}}
\newcommand{\FS}{{{\operatorname{FS}}}}
\newcommand{\supp}{{\operatorname{supp\,}}}
\renewcommand{\phi}{\varphi}
\newcommand{\eqd}{\buildrel {\operatorname{def}}\over =}
\newcommand{\go}{\mathfrak}
\newcommand{\acal}{\mathcal{A}}
\newcommand{\bcal}{\mathcal{B}}
\newcommand{\ccal}{\mathcal{C}}
\newcommand{\dcal}{\mathcal{D}}
\newcommand{\ecal}{\mathcal{E}}
\newcommand{\fcal}{\mathcal{F}}
\newcommand{\gcal}{\mathcal{G}}
\newcommand{\hcal}{\mathcal{H}}
\newcommand{\ical}{\mathcal{I}}
\newcommand{\jcal}{\mathcal{J}}
\newcommand{\kcal}{\mathcal{K}}
\newcommand{\lcal}{\mathcal{L}}
\newcommand{\mcal}{\mathcal{M}}
\newcommand{\ncal}{\mathcal{N}}
\newcommand{\ocal}{\mathcal{O}}
\newcommand{\pcal}{\mathcal{P}}
\newcommand{\qcal}{\mathcal{Q}}
\newcommand{\rcal}{\mathcal{R}}
\newcommand{\scal}{\mathcal{S}}
\newcommand{\tcal}{\mathcal{T}}
\newcommand{\ucal}{\mathcal{U}}
\newcommand{\vcal}{\mathcal{V}}
\newcommand{\zcal}{\mathcal{Z}}
\newcommand{\al}{\alpha}
\newcommand{\be}{\beta}
\newcommand{\ga}{\gamma}
\newcommand{\Ga}{\Gamma}
\newcommand{\La}{\Lambda}
\newcommand{\la}{\lambda}
\newcommand{\ep}{\varepsilon}
\newcommand{\de}{\delta}
\newcommand{\De}{\Delta}
\newcommand{\om}{\omega}
\newcommand{\Om}{\Omega}
\newcommand{\di}{\displaystyle}
%\newcommand{\l<\r>}{\left<\right>}
\newtheorem{theo}{{\sc Theorem}}
\newtheorem{maintheo}{{\sc Theorem}}
\newtheorem{cor}{{\sc Corollary}}[section]
\newtheorem{q}[theo]{{\sc Question}}
\newtheorem{p}[theo]{{\sc Problem}}
\newtheorem{conj}[theo]{{\sc Conjecture}}
\newtheorem{maincor}[maintheo]{{\sc Corollary}}
\newtheorem{maindefn}[maintheo]{{\sc Definition}}
\newtheorem{lem}[cor]{{\sc Lemma}}
\newtheorem{mainprop}[maintheo]{{\sc Proposition}}
\newtheorem{prop}[cor]{{\sc Proposition}}
\numberwithin{equation}{section}
\newenvironment{example}{\medskip\noindent{\it Example:\/} }{\medskip}
\newenvironment{rem}{\medskip\noindent{\it Remark:\/} }{\medskip}
\renewcommand{\O}[1]{O_{#1}}

\newtheorem{defn}[theo]{{\sc Definition}}
\newenvironment{claim}{\medskip\noindent{\it Claim:\/} }{\medskip}
\usepackage{hyperref}





\title[Trace estimates and improved pointwise bounds]{Trace estimates and improved pointwise bounds for joint Eigenfunctions}


\author{Xianchao Wu and Xiao Xiao}
\address{School of Mathematics and Statistics, Wuhan University of Technology, Wuhan, Hubei, China}
\email{xianchao.wu@whut.edu.cn } 
\address{Department of Mathematics and Statistics, McGill University, Montr\'eal, Quebec, Canada}
\email{xiao.xiao5@mail.mcgill.ca}


\date{}

\begin{document}

\begin{abstract}
For $L^2$-normalized joint eigenfunctions in a quantum integrable system, \cite{GT2020} gave polynomial improvements over the standard H\"omander bounds for typical points. In this paper, we improve their result by establishing a sharp bound of $h^{\frac{-n+k+1}2}$ for the points satisfying a {\em rank $k$} non-degeneracy condition. 
\end{abstract}

\maketitle


\section{Introduction}\label{section1}

Let $(M,g)$ be a compact Riemannian manifold without boundary. By basic functional analysis, the Laplace-Beltrami operator $\Delta_g:=\textit{div}_g\circ\nabla_g$ has discrete spectrum $\{\lambda_j\}_{j=0,1,...}$. The corresponding eigenfunctions are $u_j$ satisfying
$$
-\Delta_gu_j=\lambda_ju_j.
$$
We opt to use the semiclassical notation. Let $h_j:=\lambda_j^{-\frac{1}{2}}$, and omitting the subscript $j$, the above equation becomes
\begin{equation}\label{1}
    (-h^2\Delta_g-1)u_h=0.
\end{equation}
The principal symbol of the operator $P(h)=-h^2\Delta_g-1$ is 
\begin{equation}\label{2}
p(x,\xi)=|\xi|_{g(x)}^2-1.    
\end{equation}
The symbol $p(x,\xi)$ is a smooth function on $T^*M$ and can be viewed as a Hamiltonian. The associated Hamiltonian flow is exactly the lifted geodesic flow on $T^*M$. (Under the standard projection, this flow is the geodesic flow on $M$.) By the quantum-classical correspondence, (\ref{1}) and (\ref{2}) should influence each other in the following sense: The dynamical properties of the geodesic flow (\ref{2}) of $(M,g)$ influence the asymptotic properties of $u_h$ in (\ref{1}) as $h\to 0^+$.

There has been an perennial endeavor to understand the concentration of eigenfunctions. A compelling question is how spiky can a sequence of eigenfunction be: Seeking for an $L^{\infty}$ bound of $L^2$ normalized $u_h$ in terms of $h$. 

Beginning in the 1950's, in a series of works by Levitan \cite{Lev}, Avakumovi\'c \cite{Ava}, and H\"ormander \cite{Hor}, the following estimate is obtained (known as the H\"ormander bound):
\begin{equation}\label{known}
\|u_h\|_{L^\infty}=O(h^{\frac{1-n}2}),
\end{equation}
as $h\to 0$ and \eqref{known} is saturated on the round sphere. This bound was improved to $o(h^{\frac{1-n}2})$ by Safarov, Sogge, Toth, Zelditch and Galkowski \cite{Saf88, SZ02, STZ11, SZ16a, SZ16b, GT17, Gal19} under various dynamical assumptions at $x$. When $(M,g)$ has no conjugate points, a quantitative improvement
\begin{equation}\label{know1}
\|u_h\|_{L^\infty}=O(h^{\frac{1-n}2}/\sqrt{|\ln h|}),
\end{equation}
as $h\to 0$, has been known since the classical work of B\'erard \cite{Ber77, Bon17, Ran78}. In recent times, Canzani and Galkowski \cite{CG21, CG23} developed the tool of geodesic beams to study the quantitative improvements without global geometric assumptions on $(M,g)$. 

\subsection{Quantum completely integrable (QCI) system}
A semiclassical pseudodifferential operator $P_1(h)$ is called quantum completely integrable (QCI) if there exist functionally independent semiclassical pseudodifferential operators $P_2(h),\cdots,P_n(h)$ such that 
$$
[P_i(h),P_j(h)]=0,\quad i,j=1,\cdots, n.
$$
Given real numbers $E_1(h),\cdots, E_n(h)$, we want to study the $L^2$ normalised functions $u_h$ that solves $P_j(h)u_h=E_j(h)u_h$ for all $j=1,\cdots,n$. We call them joint eigenfunctions corresponding to the joint eigenvalues $E(h):=(E_1(h),\cdots,E_n(h))\in \R^n$ of the system $\hat{\mathcal{P}}:=(P_1(h),\cdots, P_n(h))$.

The principal symbols $p_j\in C^{\infty}(T^*M)$ of each operator $P_j(h)$ give rise to a moment map that defines an associated classical integrable system. The moment map is denoted by
$$
\mathcal{P}:=(p_1,\cdots,p_n):T^*M\to \R^n.
$$

The regular set of $\mathcal{P}$, denoted by $(T^*M)_{reg}$, consists of points $(x,\xi)\in T^*M$ such that
$$
\dim span\{ dp_1(x,\xi),\cdots,dp_n(x,\xi)\}=n.
$$
As a global assumption, the classical system $\mathcal{P}$ is {\textit{Liouville integrable}}, i.e. the set $(T^*M)_{reg}$ is an open dense subset of $T^*M$.

The above construction induces a Lagrangian torus fibration of $T^*M$. If $E=(E_1,\cdots,E_n)$ is a regular value of $\mathcal{P}$, then by the Arnold-Liouville theorem, $\mathcal{P}^{-1}(E)$ consists of a finite disjoint union of Lagrangian tori:
$$
\mathcal{P}^{-1}(E)=\bigcup_{j=1}^N \Lambda_j(E).
$$
Each $\Lambda_j(E)$ is a Lagrangian submanifold of $T^*M$ and is fixed by the torus action of $\mathcal{P}$. In fact, the joint flow 
$$
G_{t}(x,\xi):=\exp (t_1H_{p_1})\circ \cdots\circ\exp(t_nH_{p_n})(x,\xi)
$$
is defined for all $t\in \R^n$ and the Hamiltonian flow $\exp (t_jH_{p_j})$ preserves the level set $\{p_j=E_j\}$. Furthermore, the definition does not depend on the order of composition since we assumed $[P_i,P_j]=0$, which implies $\{p_i,p_j\}=0$.

The singular set $T^*M\setminus(T^*M)_{reg}$ can be view as the degenerate Lagrangian tori.

Galkowski-Toth \cite{GT2020} investigated the pointwise bounds of $u_h$ for the case $P_1(h)$ being a Schr\"odinger operator. More explicitly, when $E_1$ is a regular value of $p_1$ and the system is of {\em Morse type} at point $x$, they established the following bound:
\begin{equation}\label{known2}
|u_h(x)|=O(I_n(h)),
\end{equation}
here $I_2(h)=h^{-\frac{1}4}$, $I_3(h)=h^{-\frac{1}2}|\ln h|^{\frac{1}2}$ and $I_{n\geq 4}(h)=h^{\frac{2-n}2}$.

\subsection{Main results}
One purpose of this paper is to obtain a little $o$ improvement of the pointwise bounds \eqref{known2} under certain non-degeneracy condition. We achieve this in Theorem \ref{main1}. A useful tool is developed along the way (Theorem \ref{trace}), which we call a {\textit{quantitative trace estimate}}. It is of independent interest. 

We say that $p_1$ is of real principal type on the hypersurface $p_1^{-1}(E_1)$: if $E_1>0$ is a regular value of $p_1$, for any $(x,\xi)\in p_1^{-1}(E_1)$, the following inequality holds,
\begin{equation}\label{principal}
\partial_\xi p_1(x,\xi)\neq 0.
\end{equation}
One sets
\begin{equation}\label{loop set}
\mathcal{C}_x:=\{\xi\in T_x^*M:\, p_1(x,\xi)=E_1\}.
\end{equation} 
and define the {\textit{loop set}} to be
\begin{equation}
    \mathcal{L}_x:=\{\xi\in \mathcal{C}_x:\exists t<+\infty,\, \Phi_{t,p_1}(x,\xi)\in \mathcal{C}_x\},
\end{equation}
where $\Phi_{t,p_1}$ is the Hamiltonian flow of $p_1$. If furthermore, a covector loops back in a direction infinitly close to it original direction infinitly many times in the future and the past, we call it a \textit{recurrent direction}. Define the \textit{recurrent set} to be

\begin{equation}\label{recurrent set}
\mathcal{R}_x:=\{\xi\in \mathcal{C}_x:\Big(\bigcap_{T>0}\overline{\bigcup_{t>T}\Phi_{t,p_1}(x,\xi)|_{\mathcal{C}_x}}\Big)\bigcap\Big(\bigcap_{T>0}\overline{\bigcup_{t>T}\Phi_{-t,p_1}(x,\xi)|_{\mathcal{C}_x}}\Big)\}.
\end{equation}
The closure is under the subspace topology from $\mathcal{C}_x$.

It is the omega-limit of the following dynamical system: Define the return-time function $T_x:\mathcal{C}_x\to \R\cup\{\infty\}$, 
$$
T_x(\xi)=\inf_{t>0}\{t:\Phi_{t,p_1}(x,\xi)\in \mathcal{C}_x\}.
$$
Then the return map $\widetilde\Phi_x:=\Phi_{T_x,p_1}:\mathcal{L}_x\to \mathcal{C}_x$ restricted to $\widetilde{\mathcal{L}}_x:=\bigcap_{k\in \Z}\widetilde{\Phi}^k(\mathcal{L}_x)$ is a dynamical system. By (\ref{recurrent set}), $\xi\in \mathcal{R}_x$ if and only if $\xi$ is a limit point of the orbit of $\widetilde\Phi_x(\xi)$. See \cite{STZ11} and \cite{Gal19}.

Note that we have the following inclusion:
$$
\mathcal{R}_x\subset\widetilde{\mathcal{L}}_x\subset\mathcal{L}_x\subset \mathcal{C}_x.
$$
%\begin{rem}
%In the homogeneous case where $p_1=|\xi|_g^2$ and $E_1=1$. One has that \eqref{principal} is automatically satisfied. 
%\end{rem}
% Now we can explain the {\em strictly monotonic} condition.

\begin{defn}\label{def}
Suppose that $1\leq k\leq n-1$ and $\Gamma$ is a subset of $\mathcal{C}_x$. We say that $\mathcal{P}=(p_1, \dots, p_n)$ satisfies {\bf rank $k$ condition} at $x\in M$ for $\Gamma$ if
\begin{equation}\label{rk}
    \text{ for any } \xi\in\Gamma,\quad\text{dim span}\{\partial_\xi p_2 ,\,\partial_\xi p_3,\cdots, \partial_\xi p_n\} = k
\end{equation}
as a subspace in $T_{\xi}\mathcal{C}_x$, where $\partial_{\xi}$ denotes the full gradient along $\mathcal{C}_x$.

For brevity, we say $\mathcal{P}|_{\Gamma}$ is rank $k$ at $x$.
\end{defn}

\begin{rem}
    It is noteworthy that similar conditions were used in \cite[Def. 1.2]{EGK24}, under the name “fiber rank $k$ condition”, and in \cite[Thm 0.1 (2)]{Tacy19}, under the name “admissibility condition”. Our formulation is slightly more general in the sense that it specifies the energy level and the subset $\Gamma$ of energy surface of $p_1$. Geometrically, condition (\ref{rk}) means that $\mathcal{P}^{-1}(E)\cap \Gamma$ is a $k$ dimensional submanifold which can be parameterized by $(\xi_{i_1},\dots, \xi_{i_k})$.
\end{rem}

From now on, we always suppose that $\hat{\mathcal{P}}=\{\sqrt{-h^2\Delta},P_2(h),\dots, P_n(h)\}$ is a QCI system on a smooth, compact manifold without boundary, $(M,g)$. Let $\{u_h\}$ be the $L^2$-normalized joint eigenfunctions.
As a key ingredient of our paper, we can state our first theorem,
\begin{maintheo}\label{trace}
Assume $A(x, hD)\in \Psi^0(M)$ with symbol $a=\sum_{j=0}^\infty a_j h^j$. If $\mathcal{P}|_\Gamma$ is rank $k$ at $x$ with $\Gamma=\mathcal{C}_x\cap \supp a$, then there exists a positive constant $h_0$ so that for $h\in (0, h_0]$, 
\begin{equation}\label{main cor eq}
|A(x,hD) u_{h}|^2=O(h^{-n+k+1}).
\end{equation}
\end{maintheo}

\begin{rem}
Section \ref{section2} shows that Theorem \ref{trace} relies on the trace estimates, which also play a critical role in proving the local Weyl law (see \cite[Section 5]{So2014}).
\end{rem}

%Theorem \ref{trace} is a trace estimate. As an application, one can get  generalized local Weyl law,
%\begin{maintheo}
%If $A(x, hD)\in \Psi^0(M)$ with symbol $a=\sum_{j=0}^\infty a_j h^j$, there exists a positive constant $h_0$ so that for any $h\in (0, h_0]$, 
%\begin{equation}
%\big|\big|
%\end{equation}
%\end{maintheo}




% \begin{defn}\label{def}
% We say that $\mathcal P=(p_1, p_2)$ is {\bf non-critical} at $x\in M$ if 
% \begin{equation}\label{C1}
% \text{Crit} \Big(p_2\,|_{\, \mathcal{C}_x \cap \{|p_2(x,\xi)-E_2|< \epsilon\}}\Big) \,\cap\, \mathcal{R}_x=\varnothing.
% \end{equation}
% for some sufficiently small positive constant $\epsilon$.
% \end{defn}

% \begin{rem}
%     The non-critical condition is not only a condition for $x$ but also one for $u_h$ implicitly. It is because $E_1$ and $E_2$ appear in (\ref{C1}), telling us which eigenfunctions are to be considered.
% \end{rem}


\begin{defn}\label{morse}
    We say that $\mathcal{P}=(p_1,\cdots,p_n)$ is {\bf Morse} at $x\in M$ if there exists a function $f\in C^{\infty}(\R^n,\R)$, such that the principal symbol $q$ of the operator $Q:=f(P_1,\cdots,P_n)$ satisfies 
    \begin{equation}
        q\big|_{\mathcal{C}_x}(\xi) \text{ is a Morse function in } \xi. 
    \end{equation}
\end{defn}

\begin{rem}
    For surfaces of revolution, let $\hat{\mathcal{P}}$ be the standard QCI system (see Section 4). Then $\mathcal{P}$ is \textit{not} Morse at north and south poles, because $P_2$ vanishes there and $f=f(P_1)$, the symbol of whom, restricted to the energy surface of $p_1$, is a constant function.
    
    Also for the triaxial ellipsoid, the standard QCI system is \textit{not} Morse at the four umbilical points \cite{GT2020}. 
\end{rem}

Our main theorem is 

\begin{maintheo}\label{main1}
If $\mathcal{P}|_{\Gamma}$ is rank $k$ \rm{($k\geq 2$ if $n>3$)} at $x$ with $\Gamma=\mathcal{R}_x\cap \bigcap\limits_{2\leq j \leq n}\{|p_j(x,\xi)-E_j|\leq \epsilon\} $ for some $\ep>0$, and $\mathcal{P}$ is Morse, then there exists a positive constant $h_0$ such that for any $h\in (0, h_0]$,
\begin{equation}\label{main1 eq}
| u_h(x) |  = o\big(I_n(h)\big),
\end{equation}
where $I_2(h)=h^{-1/4}$, $I_3(h)=h^{-1/2}|\log h|^{1/2}$ and $I_n(h)=h^{(2-n)/2}$ if $n>3$. 

Furthermore, if $\mathcal{P}|_{\Gamma'}$ is rank $k$ at $x$ with $\Gamma'=\mathcal{C}_x\cap \bigcap\limits_{2\leq j \leq n}\{|p_j(x,\xi)-E_j|\leq \epsilon\}$ for some $\ep>0$, then there exists a positive constant $h_0'$ such that for any $h\in (0, h_0']$,
\begin{equation}\label{main1 eq2}
| u_h(x) | = O(h^{\frac{-n+k+1}2}).
\end{equation}

\end{maintheo}

\begin{rem}
   \begin{enumerate}
       \item Regarding the first conclusion (\ref{main1 eq}): The pointwise bound for a general Laplace eigenfunction (without quantum integrability assumption) possess an improvement to little $o$ if $|\mathcal{L}_x|=0$, \cite{SZ02} (In fact, only need $|\mathcal{R}_x|=0$ ,\cite{STZ11}). It is tempting to ask if the same is true for joint quasimodes in the QCI case. The answer is negative. On a generic biaxial ellipsoid, if the point $x$ is on the equator, then $|\mathcal{R}_x|=|\mathcal{L}_x|=0$. However, one can construct a quasimode of magnitude $\sim  h^{-1/4}$ attained on the equator, see section \ref{section4}. That shows some extra condition such as the rank $k$ is necessary.
       \item Regarding the second conclusion (\ref{main1 eq2}): If the conditions there are satisfied for all points $x\in M$, we restore the $L^{\infty}$ bound in \cite{Tacy19}. We note that our approach is different with Tacy's. \cite{Tacy19} is based on the result of \cite{KTZ}, however, our approach is self-contained and works point-wise. Also, the condition rank $k$ arises naturally in our proof (ref. \eqref{phi low}, \eqref{rank2 1} and \eqref{rank2 2}).
   \end{enumerate}
\end{rem}


% \begin{maintheo}\label{main1}
% Let $\{u_h\}$ be the $L^2$-normalized joint eigenfunctions of commuting operators $P_1(h)$ and $P_2(h)$ on a compact smooth Riemannian manifold $(M,g)$. Let $\big(E_1(h), E_2(h)\big)\in \textnormal{Spec} (P_1(h))\times \textnormal{Spec} (P_2(h))$, and $E_j(h)=E_j+O(h)$, $j=1,2$. Assume that
% \begin{itemize}
%     \item $p_1$ is of real principal type on $p_1^{-1}(E_1)$;
%     \item $\mathcal{P}=(p_1,p_2)$ is Morse at $x\in M$.   
%     \item $\mathcal{P}=(p_1,p_2)$ is non critical at $x\in M$. 
% \end{itemize}
% Then there exists a positive constant $h_0$ such that for any $h\in (0, h_0]$,
% \begin{equation}\label{main1 eq}
% | u_h(x) |  = o\big(I_n(h)\big),
% \end{equation}
% where $I_2(h)=h^{-1/4}$, $I_3(h)=h^{-1/2}|\log h|^{1/2}$ and $I_n(h)=h^{(2-n)/2}$ if $n>3$.
% \end{maintheo}

% \begin{rem}
%   It is shown in section \ref{example} that on $\mathbb{S}^2$, when $E_1=1,\,E_2\in[0,1)$, all points but the poles satisfy (\ref{C1}). That means, for all but the highest-weight ($E_2=1$) spherical harmonics, one has the improved bound (\ref{main1 eq}). For $E_1=E_2=1$, no point satisfies (\ref{C1}), therefore the main theorem says nothing about the highest weight spherical harmonics. 
% \end{rem}

% \begin{rem}
% The estimate \eqref{main1 eq} in Theorem \ref{main1} gives an explicit polynomial improvement over the bound \eqref{known2} and above estimate is uniform over all energy values $\{E_2\in\mathbb{R};\, (E_1, E_2)\in \mathcal{P}(T^*M)\}$. From the examples constructed in Section \ref{example}, one can see that the {\em non critical} assumption is crucial. 

% %The methodology presented in current paper could be extended to general elliptic operators and smooth curves, although such an extension need more intricate change of variables. We have chosen to state our results for the Schr\"odinger operator $P_1(h)$ and geodesic segment $\gamma$ in this work as the formulation retains its elegance and encapsulates the essential techniques.
% \end{rem}







% As one can observe in some examples \cs add refs\cs, the points at which the $L^\infty$ bound can be achieved are usually rare and special. In general, little is known for the majority of points in a manifold with respect to some sequence of eigenfunctions. 

% For example, let $u_{\lambda}$ be the highest weight spherical harmonics. Then it is clear that for any point $x$ apart from the poles $x_{\pm}$, $|u_{\lambda}(x)|\leq Ce^{-\delta(x)\lambda}$. (By legendre polynomials, or in general, by energy localization, c.f. Jeff $\&$ Toth.)

% \cs Another seemingly reasonable definition is given by comparing with the point saturating the $\| \cdot\|_{L^\infty}$. This way it is independent of an absolute criteria, such as the H\"ormander bound $h^{\frac{1-n}{2}}$. \cs

% \begin{defn}
%     Given a sequence $u_h$ satisfying $P(h)u(h)=0$ on $M$, we call a point $x\in M$ an {\bf excited} point if $$\limsup _h|u_h(x)|=\Omega (h^{\frac{1-n}{2}}).$$
%     We call a point $x$ {\bf semi-excited} if 
%     $$
%     \limsup_h |u_h(x)|=O(h^{\frac{1-n}{2}+\epsilon}),
%     $$
%     and call such supremum $\epsilon>0$ its {\bf deficiency} of excitement. 
% \end{defn}

% In other words, $u_h$ has a polynomial gain at semi-excited points and the power of $h$ it gains is its deficiency. In the above example, any $x\neq x_\pm$ is a semi-excited point of deficiency $\epsilon(u_h,x)=\frac{n-1}{2}$ as they lose all the powers of $h$.

% \begin{p} Charaterize the semi-excited points and their deficiencies by the geometry and dymanics of the manifold.
% \end{p}

% \begin{conj}
%     The set of semi-excited points are of full-measure for any $u_h$.
% \end{conj}

% \begin{conj}
%     Fixing a sequence $u_h$, the function $\epsilon(u_h,\cdot)$ is lower semi-continuous. Moreover, on $M\setminus\{\text{excited points}\}$, $\epsilon$ is locally constant.
% \end{conj}

% \begin{rem}
%     Any sequence $u_h$ associates to a defect measure $\mu$. The projection of its support on the base is the singular support $SS(u_h)$. We know from energy concentration that $u_h$ decays exponentially outside of $SS(u_h)$, therefore its deficiency of excitement is $\epsilon=+\infty$. Thus, what matters is the behavior within $SS(u_h)$.
% \end{rem}

% The example of sphere case: Let $u_h$ be highest-weight spherical harmonics. Since the norm satisfies $\|(x_1+ix_2)^l\|_{L^2(S^2)}^2=\frac{\Gamma(l+\frac{1}{2})}{\Gamma(l+1)}\sim l^{-1/2}$, the $L^2$ normalized highest-weight spherical harmonics are $l^{\frac{1}{4}}(x+iy)^l$, achieving its peak value $l^{1/4}$ on the equator $\{x^2+y^2=1,\,z=0\}$. In dimension $n$, they have height $\lambda^{\frac{n-1}{4}}$. Hence the deficiency of a point $x$ on the equator is $\epsilon(u_h,x)=\frac{1-n}{4}$, and for any other point, its deficiency is $\epsilon(u_h,y)=+\infty$.

% Some known inequalities, regarding "exceptional sequences": 

% [Rudnick-Sarnak][Donnelly][Milicevic]

% Then bring up the rotational surface case, and state the main theorem.



% Now let us turn our attention to the joint eigenfunctions of quantum completely integrable systems. We begin by considering a self-adjoint semiclassical pseudodifferential operator $P_1(h): C^\infty(M)\to C^\infty(M)$ of order $m$ in the classical sense. This means that the symbol $p_1(x,\xi)$ of $P_1(h)$ satisfies the inequality $|p_1(x,\xi)|\geq c|\xi|^m-C$. We say that $P_1(h)$ is quantum completely integrable (QCI) if there exist functionally independent $h$-pseudodifferential operators $P_2(h), \dots, P_n(h)$ with property that
% \begin{equation*}
% [P_i(h), P_j(h)]=0; \quad i,j=1,\dots, n.
% \end{equation*}
% Let $E(h)=(E_1(h), \dots, E_n(h))\in \mathbb{R}^n$ denote the joint eigenvalues of the operators $P_1(h), \dots, P_n(h)$. We denote the $L^2$-normalized joint eigenfunction with joint eigenvalue $E(h)$ by $u_{E,h}$. When the joint energy value $E$ is clear from the context, we will simplify the notation and write $u_h$ instead of $u_{E,h}$.

% Galkowski-Toth \cite{GT2020} investigated the pointwise bounds of $u_h$ for Schr\"odinger operator $P_1(h)$. More explicitly, in the n-dimensional case and when $E_1$ is a regular value of $p_1$ and the system is of {\em Morse type} at point $x$, they established the following bound:
% \begin{equation}\label{known2}
% |u_h(x)|=O(I_n(h)),
% \end{equation}
% here $I_2(h)=h^{-\frac{1}4}$, $I_3(h)=h^{-\frac{1}2}|\ln h|^{\frac{1}2}$ and $I_{n\geq 4}(h)=h^{\frac{2-n}2}$.

% In this note, we aim to show an improvement of the pointwise bounds \eqref{known2} under a {\em non critical} condition. In fact, it can be observed that the highest-weight spherical harmonics at the equator saturate the upper bound \eqref{known2} (refer to the examples in Section \ref{example}). Therefore, additional assumptions regarding the point are necessary to improve the bound \eqref{known2}.





%Beginning in the 1950's, the works of Levitan, Avakumovi\'c, and H\"ormander\cite{} proved the estimate
%\begin{equation}\label{known}
%\|u_h\|_{L^\infty}=O(h^{\frac{1-n}2}),
%\end{equation}
%as $h\to 0$ and \eqref{known} is saturated on the round sphere. This bound was improved to $o(h^{\frac{1-n}2})$ by Safarov, Sogge, Toth, Zelditch and Galkowski \cite{} under various dynamical assumptions at $x$. \cite{} were the first to study $L^\infty$ bounds under local dynamical assumptions. When $(M,g)$ has no conjugate points, a quantitative improvement
%\begin{equation}\label{know1}
%\|u_h\|_{L^\infty}=O(h^{\frac{1-n}2}/\sqrt{|\log h|}),
%\end{equation}
%as $h\to 0$, has been known since the classical work of B\'erard \cite{}. And recently Canzani and Galkowski \cite{} developed the tool of geodesic beams to study the quantitative improvements without global geometric assumptions on $(M,g)$. In this paper, we aim to use a different dynamical assumption to study the $o(1)$ improvement of \eqref{known}.


%Set $P(h):C^\infty(M)\to C^\infty(M)$ with principle symbol $p(x,\xi)$ to be a self-adjoint semiclassical pseudodifferential operator of order $m$ that is elliptic in the classical sense, i.e. $|p(x,\xi)|\geq c|\xi|^m-C$. Given $E$ a regular value of $p$, we assume that for $(x,\xi)\in p^{-1}(E)$,
%\begin{equation}\label{regular}
%\partial_\xi p(x,\xi)\neq 0.
%\end{equation}
%In the homogeneous case where $p=|\xi|^2_g$ and $E=1$, by homogeneity, $\xi\cdot \partial_\xi|\xi|_g^2=2|\xi|_g^2$ so that \eqref{regular} is automatically satisfied. We set $p_x^{-1}(E):=\{\xi\in T_x^*M; p(x,\xi)=E\}$.


%\begin{rem}
%Geometrically condition \eqref{C1} tells that the curve $\tilde \gamma $ intersects $T^*\gamma$ transversely. 

%Analitically, if one assumes that $\gamma$ is
%\begin{equation*}
%\{(\tau,0): -1/2\leq \tau\leq 1/2\},
%\end{equation*} 
%due to $(E_1, E_2)\in B_{reg}$ one can apply the implicit function theory to parametrize the curve $\tilde \gamma$ by variable $\tau$. Then condition \eqref{C1} guarantees that
%\begin{equation}\label{C2}
%\big(\frac{d\xi_1}{d\tau}\big)^2+\big(\frac{d\xi_2}{d\tau}\big)^2\neq 0,
%\end{equation}
%here $\xi=(\xi_1,\xi_2)$.

%If $P_1(h)$ is Laplacian with $E_1=1$ and one constructs a polar coordinate along $\gamma$, then $\tilde \gamma=\{p_2(\tau,\cos\theta,\sin\theta)=E_2\}$. From \eqref{C1} one has that
%\begin{equation}\label{C3}
%|\partial_\theta p_2(\tau, \cos\theta, \sin\theta)|>C>0, \quad \text{when $\theta$ is closed to $\frac{\pi}2$ or $\frac{3\pi}2$}.
%\end{equation}
%In this case, \eqref{C2} is equivalent to say that
%\begin{equation}\label{C4}
%|\frac{\partial \theta}{\partial \tau}|>C'>0.
%\end{equation}
%\end{rem}


% \subsection{The rank k condition on recurrent directions}
% There are three conditions in Theorem \ref{main1}: The first one allows a slight generalization of the Laplacian and is used by many authors; The second one is technical but holds generically \cite{GT2020}. We care more about the last one.

% We wish to give a more straightforward criterion of the condition \ref{C1} in terms of the local geometry and dynamics. In fact, we have the following characterization.

% \begin{lem}\label{equiv}
%    Let $p_1,p_2$ be given as above, and $x\in M$ fixed. The following are equivalent:
%    \begin{enumerate}
%        \item $\mathcal{P}=(p_1,p_2)$ is non critical at $x$.
%        \item For any $\xi\in\mathcal{R}_x$, the two phase-space flows $\Phi_{t,p_1}$ and $\Phi_{t,p_2}$ are transversal at $(x,\xi)$.
%        \item For any $\xi\in\mathcal{R}_x$, $H_ {p_2}\notin \text{span}\{H_{p_1},H_{q_1}\}$ as vectors in $T_{(x,\xi)}T^*M$, where $q_1$ is the dual variable to $p_1$ in local Arnold-Liouville coordinates.
%    \end{enumerate}
% \end{lem}

% \begin{proof}
%     The equivalence of (2) and (3) follows from the local Arnold-Liouville theorem and the fact that $\{p_1,p_2\}=0$.

%     (1)$\implies$(3). Note that (\ref{C1}) means that for any $\xi\in \mathcal{R}_x$, $\partial(p_2|_{\{p_1=E_1\}})\neq 0$. By the implicit function theorem, that is true if and only $\{p_1=E_1\}\cap \{p_2=E_2\}$ is locally a smooth codimension-$2$ submanifold of $T^*M$ near $(x,\xi)$, $\xi\in \mathcal{R}_x$. It means their corresponding Hamiltonian vectors are not collinear. On the other hand, $\{p_1,p_2\}=0$ so $p_2$ cannot be the dual of $p_1$. Therefore, $H_{p_2}$ must be independent to $H_{p_1}$ and $H_{q_1}$.

%     (3)$\implies$(1). Assume on the contrary, that $H_{p_2}=cH_{p_1}+c'H_{q_1}$. Because $p_2$ poisson- commutes with $p_1$,
%     $$
%     0=H_{\{p_2,p_1\}}=[H_{p_2},H_{p_1}]=c'[H_{q_1},H_{p_1}]=c'H_{\{q_1,p_1\}}
%     $$
%     therefore $c'=0$, and $H_{p_2}=cH_{p_1}$. But it contradicts (\ref{C1}) by the implicit function theorem. 
% \end{proof}


% \begin{cor}
% If the non-critical condition in Theorem \ref{main1} is replaced by (2) or (3) in Lemma \ref{equiv}, then the $L^2$-normalized joint eigenfunctions $u_h$ satisfy the upper bound
% \begin{equation}\label{little o}
% |u_h(x)|\leq o\big(I_n(h)\big),
% \end{equation}
% where $I_2(h)=h^{-1/4}$, $I_3(h)=h^{-1/2}|\log h|^{1/2}$ and $I_n(h)=h^{(2-n)/2}$ if $n>3$.
% \end{cor}

% The following example of irrational translation on flat tori demonstrates that the rank $k$ condition on the set of recurrent directions plays an important role.

% \begin{example}
% Let $M=\R^n/\Z^n,\,(n>3)$  be the standard flat torus, and let $p_j=\sum_i c_i^jdx_i$ be unit covectors for $j=2,\cdots,n$, such that $\text{rank}_{\Q}(c_i^j)=k$. Without loss of generality, set $x=0$, and then $\mathcal{R}_x$ consists of all the unit vectors in rational directions. The geodesic flow on the flat torus is
% $$
% \Phi_{t,p_1}:(x,\xi)\mapsto (x+t\xi,\xi) ,
% $$
% the Hamiltonian flow of $p_2$ is an irrational translation 
% $$
% \Phi_{t,p_2}:(x,\xi)\mapsto(x+t\eta,\xi).
% $$ 
% Therefore, the two flows are transversal at $(x,\xi)$ if $\xi$ is rational, veryfying (2) of Lemma \ref{equiv}. Let us consider the joint eigenfunctions in $E_2=0$ case. ($0$ is in the spectrum of any translation operator on the torus.) We look for solutions $u_h(x_1,\cdots,x_n)$ to the system
% \begin{equation}\label{system}
% \begin{cases}
%     (-h^2\Delta-1)u_h=0\\
%     u_h(x_1,x_2,\bar{x})=u_h(x_1+c_1t,x_2+c_2t,\bar{x}),\quad \forall t\in\R.
% \end{cases}
% \end{equation}
% Since the trajectories of the irrational translation is dense in the fundamental domain, the second line of (\ref{system}) only means that $u_h$ is constant with respect to $x_1,x_2$. Therefore, using separation of variables and a simple coordinate change, the above system is reduced to the usual eigenvalue problem on $\R^{n-2}$ with periodic boundary condition. The H\"ormander bound (\ref{known}) give us $|u_h|\leq Ch^{\frac{1-(n-2)}{2}}=Ch^{\frac{3-n}{2}}$ which is indeed $o(h^{\frac{2-n}{2}})$ as in (\ref{little o}).

% For other values of $E_2$ the issue is similar. One has to average out the value of $u_h(x)$, $x\in\R^n$ in each translation of the fundamental domain $[0,1]^n$. That forces $u_h$ to be $0$ along the $x_1x_2$-plane, reducing the dimension of the system by $2$.
% \end{example}

% \begin{rem}
% For a general $\eta=(\eta_1,\cdots,\eta_n)$, write $k=\dim_{\Q}\{\eta_1,\cdots,\eta_n\}$. The flow line of $\eta$ foliates a $k$-dimensional subspace $W_k$ of $\R^n$ and problem is reduced to a system like (\ref{system}) replacing the second line by $u_h|_{W_k}=const.$ Then the H\"ormander bound shows $|u_h|\leq Ch^{\frac{k+1-n}{2}}$. Without further assumptions (e.g. one of the equivalent conditions in Lemma \ref{equiv}), $k$ might be $1$ and we do not have the desired little $o$ improvement. The non critical condition implies $k\geq 2$, and we get at least as good as $h^{\frac{3-n}{2}}$. Similar implications of irrationality on tori was observed in \cite{TZ02}.
% \end{rem}


%\subsection{Outline of the paper}


%\noindent{\sc Acknowledgements:} 

% \subsection{1-dimensional Symmetries}

% Symmetry on the manifold plays an important role in the study of norm bounds of eigenfunctions, e.g. \cite{Letter}, \cite{M16}, on rank $r$ symmetric spaces, \cite{TZ02},\cite{TZ03},\cite{GT2020} on QCI systems, and \cite{Tacy19},\cite{Tacy21} on joint quasimodes.

% % Take the round sphere $\mathbb{S}^n$ as a simple example. It is known that the zonal spherical harmonics on the round sphere saturate the bound (\ref{known}). Thus, one cannot expect to improve on the global results by just presuming a $1$-dimensional symmetry. However, such optimality only occurs for special eigenfunctions (zonal) at special points (poles). Therefore, a natural question is to find conditions on $u_h$ and/or conditions on $x$ such that $|u_h(x)|=o(h^{\frac{1-n}{2}})$ or even better. We found such a condition, definition \ref{def}, that makes use of the $1$-dimensional symmetry of the manifold in an essential way.

% In this work, we care about the following manifolds (in increasing order of generalities):  
% $$
% \begin{Bmatrix}
%     \text{Surfaces}\\
%     \text{of revolution}
% \end{Bmatrix}
% \subset
% \begin{Bmatrix}
%     \text{Riemannian}\\
%     \text{warped products}
% \end{Bmatrix}
% \subset
% \begin{Bmatrix}
%     \text{Riemmanian manifolds}\\
%     \text{with a Killing field}
% \end{Bmatrix}
% $$
% In those cases, there is a natural choice of a second Hamiltonian $p_2$ that Poisson-commutes with the Hamiltonian for the geodesic flow, $p_1=|\xi|^2_g$. 

% Let $(M,g)$ be a Riemannian manifold with a Killing vector field $X=\sum_{j=1}^nc_j(x)\partial_{x_j}$ in some local coordinate chart. We view $-ihX:C^{\infty}(M)\to C^{\infty}(M)$ as an $h$-pseudodifferential operator of zeroth order, and set $P_2(h)=-ihX$. Then $p_2=\sum_{j=1}^nc_j(x)\xi_j$ in the natural local chart in $T^*M$. 

% It is well known that a vector field $X$ satisfies $[X,\Delta_g]=0$ if and only if $X$ is Killing. In fact, for any vector field we have $[X,\Delta_g]=\frac{d}{dt}\Big\vert_{t=0}\Phi_t^*\Delta_g$. But $X$ is Killing if and only if its flow $\Phi_t$ is a smooth family of isometries for small $t$, along which the Laplacian is invariant.

% We can compute the Hamiltonian vector field $H_{p_2}$ in the canonical local coordinates of $T^*M$. It turns out to be the cotangent lifting of $X$.
% \begin{equation}\label{Hp2}
% H_{p_2}=-\sum_{j,k=1}^n\partial_{x_k}c_j(x)\xi_j\partial_{\xi_k}+\sum_{j=1}^nc_j(x)\partial_{x_j}.
% \end{equation}
% It is straightforward to check $\iota_{H_{p_2}}\omega=\sum_{j,k}\partial_{x_k}c_j(x)\xi_jdx_k+\sum_j c_j(x)d\xi_j=dp_2$. On the other hand,
% \begin{equation}\label{Hp1}
%     H_{p_1}=-\sum_{j,k,l=1}^n\partial_{x_k}g^{jl}(x)\xi_j\xi_l\partial_{\xi_k}+2\sum_{j,k=1}^ng^{jk}(x)\xi_k\partial_{x_j}
% \end{equation}


% % \begin{lem}
% %    If $X$ is recurrent at $x$, i.e. $X^{\flat}\in \mathcal{R}_x$, and for any $\xi$ such that $\xi^{\sharp}\parallel X$, we have
% %    \begin{equation}
% %        \text{rank}
% %        \begin{pmatrix}
% %            \sum_j\partial_{x_1}c_j(x)\xi_j&\cdots& \sum_j\partial_{x_n}c_j(x)\xi_j\\
% %            \sum_{j,l}\partial_{x_1}g^{jl}(x)\xi_j\xi_l&\cdots&\sum_{j,l}\partial_{x_n}g^{jl}(x)\xi_j\xi_l
% %        \end{pmatrix}_{n\times 2}<2.
% %    \end{equation}
% % Then, $\mathcal{P}=(p_1,p_2)$ is critical at $x$. The converse is also true.
% % \end{lem}

% \begin{lem}\label{dichotomy}
%     Assume $X(x)\neq 0$. $\mathcal{P}=(p_1,p_2)$ is non-critical at $x$ if and only if either of the following holds on $\{(x,\xi):|\sum_{j=1}^nc_j(x)\xi_j-E|<\epsilon\}$ for some $\epsilon>0$:
%     \begin{enumerate}
%         \item $X$ is not recurrent, i.e. $\frac{X^{\flat}}{|X^{\flat}|}\notin \mathcal{R}_x$;
%         \item $X$ is recurrent, but $\exists k\in\{1,...,n\}$ such that
%         \begin{equation}\label{else}
%             \frac{\sum_j\partial_{x_k}c_j(x)\xi_j}{\sum_{j,l}\partial_{x_k}g^{jl}(x)\xi_j\xi_l}\neq |X|\sgn(X).
%         \end{equation}
%         Where we define $\sgn(X)=1 $ if $\langle X^{\flat},\xi\rangle>0$, and $\sgn(X)=-1 $ if $\langle X^{\flat},\xi\rangle<0$.
%     \end{enumerate}
% \end{lem}
% \begin{proof}
%     Use (3) of Lemma \ref{equiv}. From the proof of Lemma \ref{equiv}, $[P_2,\Delta_g]=0\implies H_{p_2}$ has no $H_{q_1}$ components. Therefore, we only need to check if $H_{p_2}$ and $H_{p_1}$ are collinear. 
   
%     By the explicit formulas (\ref{Hp2}) and (\ref{Hp1}), it is sufficient that 
%     \begin{align}\label{collinear}
%     \sum_{j=1}^nc_j(x)\partial_{x_j}\quad\text{and}\quad 2\sum_{j,k=1}^ng^{jk}(x)\xi_k\partial_{x_j}\quad\text{not collinear for any }  \xi\in\mathcal{R}_x.
%     \end{align}
%     Since $X\neq 0$ at $x$, (\ref{collinear}) means that $X$ and $\xi^{\sharp}$ not collinear for any $\xi\in \mathcal{R}_x$, which is equivalent to that $X^{\flat}$ does not point to a recurrent direction.

%     If this good scenario does not occur, that is, if $\frac{X^{\flat}}{|X^{\flat}|}$ is unluckily a recurrent direction, then $X^{\flat}=\langle X^{\flat},\xi\rangle\xi=\sgn(X)|X^\flat|\xi=\sgn(X)|X|\xi$ for some $\xi\in \mathcal{C}_x$. For $H_{p_2}$ and $H_{p_1}$ to be independent, we need some of the coefficient ratio in (\ref{Hp2}), (\ref{Hp1}) to be different from $\sgn(X)|X^{\flat}|$, leading to the condition (\ref{else}).
% \end{proof}

% Using the above, we can state our second main result.

% \begin{maintheo}
% Let $(M,g)$ be a compact Riemannian manifold with dimension at least $2$, and let $X$ be a Killing vector field, $E(h)\in \textnormal{Spec}(-ihX)$, $E(h)=E+O(h)$. If $X$ satisfies either (1) or (2) in Lemma \ref{dichotomy} on $\{(x,\xi):|\sum_{j=1}^nc_j(x)\xi_j-E|<\epsilon\}$ for some $\epsilon>0$, and assume $X\neq 0$ at $x$. Then the $L^2$-normalized solutions $u_h$ to the system
% \begin{equation}
%     \begin{cases}
%         (-h^2\Delta_g-1)u_h=0\\
%         -ihXu_h=E(h)u_h
%     \end{cases}
% \end{equation}
% satisfy the upper bound
% \begin{equation}
% |u_h(x)|\leq o\big(I_n(h)\big),
% \end{equation}
% where $I_2(h)=h^{-1/4}$, $I_3(h)=h^{-1/2}|\log h|^{1/2}$ and $I_n(h)=h^{(2-n)/2}$ if $n>3$.
% \end{maintheo}
% \begin{proof}[Idea of proof]
%     The Crucial points are:
%     \begin{enumerate}
%         \item[1.] To show that the Morse condition at $x$ translates to $X(x)\neq 0$.
%         \item[2.] To show that the $[h^2\Delta, hX]=0$ translates to $\{|\xi|^2,\sum c_j\xi_j\}=O(h)$, which is not strictly $0$. But we can still make use of the conditions in Lemma \ref{equiv} or Lemma \ref{dichotomy}, because they are open conditions.
%     \end{enumerate}
    
% \end{proof}
% Then, we can make use of the techniques in treating the QCI case. Our novelty is that, in this case $u_h$ is supposed to be the joint eigenfunciton of only two operators in an $n$-dimensional manifold.

% Let $P_1(h)=-h^2\Delta_g$ be the Laplacian, and $P_2(h)$ be an $h$-pseudodifferential operator with
% $$
% [P_1(h),P_2(h)]=0.
% $$
% Let $E(h)=(E_1(h),E_2(h))\in \R^2$ denote the joint eigenvalue of the operators $P_1(h),P_2(h)$. We denote the $L^2$-normalized joint eigenfunction with joint eigenvalue $E(h)$ by $u_{E,h}$, or simply $u_h$ if $E$ is clear from the context.

% For instance, let $M$ be $\mathbb{S}^2$ with round metric and $P_2=-ih\partial_{\phi}$. The joint eigenfunctions for $E(h)=(1,0)$ and $E'(h)=(1,1)$ are solutions to the systems
% \begin{equation}
%     \begin{cases}
%         -h^2\Delta u_h=u_h\\
%         -ih\partial_{\phi} u_h=0
%     \end{cases}
%     \quad
%     \text{and}
%     \quad
%     \begin{cases}
%         -h^2\Delta u_h=u_h\\
%         -ih\partial_{\phi} u_h=u_h.
%     \end{cases}
% \end{equation}
% respectively. By separation of variables, one can solve these systems and get zonal and highest-weight spherical harmonics respectively.

\subsection{Example and superintegrability}\label{section superint}
% Consider a biaxial ellipsoid 
% $$
% M=\Big\{(x_1,x_2,x_3)\in \R^3:\frac{x_1^2}{a^2}+\frac{x_2^2}{a^2}+\frac{x_3^2}{c^2}=1\Big\},
% $$
% $$
% p_1=\xi_t^2+f^{-2}(t)\xi_{\phi}^2,\quad p_2=\xi_{\phi}.
% $$
% \cs We have that $\mathcal{P}|_{\{\xi_{t}\neq 0\}}$ satisfies rank $1$ condition, that is, as long as the covector $\xi$ is not horizontal. If constants $a,c, E_1, E_2$ are well chosen (see section \ref{section4}) and $x$ is not a point on the equator, one can arrange that horizontal covectors are not recurrent, i.e. $\mathcal{R}_x\subset\{\xi_t\neq0\}$. Therefore, the first conclusion in theorem \ref{main1} holds. 

% To see that the rank $k$ condition cannot be removed, we simply consider the highest-weight spherical harmonics. It saturates the $u_h(x)=Ch^{-1/2}$ bound for $x$ on the equator, and decays exponentially for any $x$ outside the equator and the poles.

% In the last section, we construct a quasimode near the "equator" of a surface of revolution, which has the same behavior as the highest-weight spherical harmonics. Such a quasimode has $h^{-1/4}$ sup norm and attains it at $x$ on the "equator", making (\ref{main1 eq}) impossible for $n=2$.

Let $M=\mathbb{T}^n$, $n\geq 3$, then $\mathcal{C}_x$ is the round cosphere. Let $p_2=\xi_2,\cdots,p_n=\xi_n$. On any pair of poles, at most one of $p_j$'s is degenerate. Therefore, $\mathcal{P}|_{\mathcal{C}_x}$ only satisfies a rank $n-2$ condition at any $ x\in \mathbb{T}^n$. As a consequence, (\ref{main1 eq2}) gives $|u_h(x)|=O(h^{-1/2})$, a bound that is not optimal. However, we can use a simple trick as follows:

Let $u_h$ be the joint eigenfunction of 
$$\hat{\mathcal{P}}=\{\sqrt{-h^2\Delta},-ih\partial_{x_2},-ih\partial_{x_3},\cdots,-ih\partial_{x_n}\}$$
with eigenvalue $E=(1,1,0,\cdots,0)$, then $\mathcal{P}$ degenerates at $(|\xi|,\xi_2,\cdots,\xi_n)=(1,1,0,\cdots ,0)$. But the same function $u_h$ can be viewed as a joint function 
$$
\hat{\mathcal{Q}}=\{\sqrt{-h^2\Delta},-ih\partial_{x_1},-ih\partial_{x_3},\cdots,-ih\partial_{x_n}\}
$$
with eigenvalue $F=(1,0,0,\cdots,0)$. $\mathcal{Q}$ satisfies the rank $n-1$ condition on $\Gamma'$. As a consequence, (\ref{main1 eq2}) gives $|u_h(x)|=O(1)$ which is sharp.

\begin{figure}[h]
\includegraphics[width=0.9\textwidth]{drawing1.pdf}
\caption{On the left: At $\xi$, the system $\mathcal{P}$ only satisfies a rank $n-2$ condition; On the right: At $\xi$, the system $\mathcal{Q}$ satisfies the rank $n-1$ condition. The same function $u_h$ is a joint eigenfunction of both $\hat{\mathcal{P}}$ and $\hat{\mathcal{Q}}$.}
\label{Figure 1}
\end{figure}

This method is possible only if there are at least $n$ independent, commuting operators that also commute with the Laplacian. We will not delve deeper into this issue in this paper.



\subsection{Plan of the paper} Section \ref{section2} proves Theorem \ref{trace} by a Fourier inversion of the operator and computing the on-diagonal Schwartz kernel, which we call the trace estimate (\`a la Sogge). It is an important technical ingredient of this paper. In section \ref{section3}, we finish the proof of Theorem \ref{main1} by taking a closer look at the recurrent set and building test symbols accordingly. We conclude by a study of examples on surfaces of revolutions in section \ref{section4}.

\subsection*{Acknowledgments} We thank Weiwei Wang for valuable discussions at the beginning of preparing this paper.


\section{Trace estimates}\label{section2}

In this section, we aim to prove Theorem \ref{trace}. Firstly we establish following trace estimate which is similar to Proposition 2.2 in \cite{SZ02}, but in QCI case.

\begin{lem}\label{lwl}
Assume $A(x, hD)\in \Psi^0(M)$ with symbol $a=\sum_{j=0}^\infty a_j h^j$. Let $\rho\in C_0^\infty(\mathbb{R})$ vanish for $|t|\geq \epsilon_0$, $\epsilon_0>0$ small enough, and satisfy $\rho(0)=1$. Let $\Gamma:=\mathcal{C}_x\cap \supp a$. Then if $(p_1|_{\Gamma}, p_2|_{\Gamma})$ is {\bf rank $1$} at $x$, one has that
\begin{equation}\label{main lemma eq2}
\int_{\mathbb{R}} \hat{\rho}(t) \big(A e^{\frac{i}h t(P_1-E_1)}  \rho\big(h^{-1}[P_2(h)-E_2]\big)  A^*\big)(t,x,x) dt =O(h^{-n+2}).
\end{equation}

Otherwise, one has the following
\begin{equation}\label{main lemma eq}
\int_{\mathbb{R}} \hat{\rho}(t) \big(A e^{\frac{i}h t(P_1-E_1)}  \rho\big(h^{-1}[P_2(h)-E_2]\big)  A^*\big)(t,x,x) dt=O(1) I_h(x)^2,
\end{equation}
if $(p_1, p_2)$ is Morse type.
\end{lem}

\begin{proof}

We can write the Schwartz kernel of $e^{\frac{i}h t(P_1-E_1)}$ in the form 
\begin{equation}
(2\pi h)^{-n} \int_{\mathbb{R}^n}e^{\frac{i}h [\varphi_1(t, y_1,\xi_1)-\left<w_1, \xi_1\right>-tE_1]} a_1(t, y_1,w_1,\xi_1;h) d\xi_1+ O_{C^\infty}(h^\infty)
\end{equation}
where $a_1\sim \sum_{j=0}^\infty a_{1,j} h^j$, $a_{1,j}\in C^\infty$, $a_{1,0}\geq C>0$ and $\varphi_1(t, y_1,\xi_1)$ solves the eikonal equation
\begin{equation}\label{phi1}
\partial_t \varphi_1(t, y_1,\xi_1)=p_1(y_1, \partial_{y_1} \phi_1), \quad \varphi_1(0, y_1, \xi_1)=\left<y_1,\xi_1 \right>.
\end{equation}
And
\begin{equation}
\rho\big(h^{-1}[P_2(h)-E_2]\big)=\frac{1}{2\pi}\int_{\mathbb{R}}\hat{\rho}(s) e^{\frac{i}h s(P_2(h)-E_2)} ds.
\end{equation}
Here $e^{\frac{i}h s(P_2(h)-E_2)}$ has a Schwartz kernel that is of the form
\begin{equation}
(2\pi h)^{-n} \int_{\mathbb{R}^n}e^{\frac{i}h [\varphi_2(s, w_1,\xi_2)-\left<y_2, \xi_2\right>-sE_2]} a_2(s, w_1,y_2,\xi_2;h) d\xi_2+ O_{C^\infty}(h^\infty),
\end{equation}
where $a_2\sim \sum_{j=0}^\infty a_{2,j} h^j$, $a_{2,j}\in C^\infty$, $a_{2,0}\geq C>0$ and $\varphi_2(s,w_1,\xi_2)$ solves the eikonal equation
\begin{equation}\label{phi2}
\partial_s \varphi_2(s, w_1, \xi_2)=p_2(w_1, \partial_{w_1} \phi_2), \quad \varphi_2(0, w_1, \xi_2)=\left<w_1,\xi_2 \right>.
\end{equation}
From \eqref{phi1} and \eqref{phi2}, one can derive the following Taylor expansions for $\varphi_1(t, y_1,\xi_1)$ (resp. $\varphi_2(s, w_1,\xi_2)$) centered at $t=0$ (resp. $s=0$).
\begin{align}
\varphi_1(t, y_1,\xi_1)&=y_1\xi_1+ tp_1(y_1,\xi_1)+O_{y_1,\xi_1}(t^2), \\
\varphi_2(s, w_1,\xi_2)&=w_1\xi_2+ sp_2(w_1,\xi_2)+O_{w_1,\xi_2}(s^2).
\end{align}

%Hence \eqref{main kernel1} is equal to
%\begin{equation}\label{main kernel}
%h^{- 4}\int \hat{\rho}(s) \hat{\rho}(t) \exp[\frac{i}h \Phi(s,t,x,w,y,\xi, \eta)]  
%c'(s,t,x,w,y,\xi, \eta;h) dw d\xi dt d\eta ds + O(h^\infty).
%\end{equation}


%One can apply stationary phase in $(w, \xi)$-variables in \eqref{main kernel}. The critical point equations for $(w, \xi)$ are
%\begin{align}
%\xi&=\eta+ t\partial_w p_1(w,\xi)+t^2\partial_w f+ O(t^3), \\
%w&=y+t\partial_\xi p_1(w,\xi)+t^2\partial_\xi f+O(t^3).
%\end{align}
%With applying Taylor expansion at the critical point $(w_c, \xi_c)$ one has that \eqref{main kernel} equals
%\begin{align*}
%&h^{-2}\int \hat{\rho}(s) \hat{\rho}(t)  \exp\Big[\frac{i}h\Big(\left<x-y,\eta\right>+ t(p_1(y,\eta)-E_1\big)+s\big(p_2(x,\eta)-E_2\big)  \nonumber \\
%&+t^2 f(y, \eta) +s^2 g(x,\eta)+O_{x,\eta}(s^3)+O_{y,\eta}(t^3) \Big)\Big] c_1(s,t,x,y, \eta;h) d\eta dt ds =: \tilde{I} (x,y; h),
%\end{align*}
%here $c_1\sim \sum_{j=0}^\infty c_{1,j} h^j$ and $c_{1,j}\in C^\infty$.

%Now set
%\begin{align}
%\tilde{I} (x,y; h):=
%h^{-2}\int \hat{\rho}(s)  \hat{\rho}(t)  &\exp[\frac{i}h\Big(\left<x-y,\eta\right>+ t(p_1(y,\eta)-E_1\big)+s\big(p_2(x,\eta)-E_2\big)  \nonumber \\
%&+O_{x,\eta}(s^2)+O_{y,\eta}(t^2) \Big)] c() d\eta dt ds
%\end{align}

In conclusion, we need to prove that
\begin{align}\label{main kernel}
h^{- 4n}\int &\hat{\rho}(s)\hat{\rho}(t) \exp[\frac{i}h \left<x-y_1, \xi\right>] \exp[\frac{i}h \Phi(s,t,y_1, w_1, y_2, \xi_1, \xi_2)]\exp[\frac{i}h \left<y_2-x, \eta\right>]\cdot  \nonumber \\
& c(s,t,y_1,w_1,y_2,\xi_1,\xi_2;h) a(x, \xi) \overline{a(x, \eta)} dy_1 d\xi dw_1 d\xi_1 d\xi_2 dy_2 d\eta dt ds,
\end{align}
can be bounded by the right hand side of \eqref{main lemma eq2} or \eqref{main lemma eq}, here $c(\cdot)=a_1(t, y_1,w_1,\xi_1;h) a_2(s,w_1,y_2,\xi_2;h) $ and
\begin{align*}
\Phi(s,t,y_1,w_1,y_2,\xi_1, \xi_2)=&\left<y_1-w_1,\xi_1 \right>+\left<w_1-y_2, \xi_2\right> +t\big(p_1(y_1,\xi_1)-E_1\big)+s\big( p_2(w_1,\xi_2)-E_2 \big)\nonumber \\
&+O_{y_1,\xi_1}(t^2)+O_{w_1,\xi_2}(s^2).
\end{align*}

Now we set
\begin{equation*}
c'(s,t,y_1,w_1,y_2,\xi_1,\xi_2;h)=\chi\big(p_1(y_1,\xi_1)-E_1\big) \chi\big(p_2(w_1,\xi_2)-E_2\big)c(s,t,y_1,w_1,y_2,\xi_1,\xi_2;h),
\end{equation*}
here $\chi\in C_0^\infty(\mathbb{R})$ vanish for $|t|\geq \epsilon$, $\epsilon>0$ small enough, and satisfy $\chi(0)=1$.
Hence \eqref{main kernel} is equal to
\begin{align}\label{main kernel1}
h^{- 4n}\int &\hat{\rho}(s) \hat{\rho}(t) \exp[\frac{i}h \left<x-y_1, \xi\right>] \exp[\frac{i}h \Phi(s,t,y_1,w_1,y_2,\xi_1,\xi_2)]\exp[\frac{i}h \left<y_2-x, \eta\right>]  \nonumber \\
 &c'(s,t,y_1,w_1,y_2,\xi_1,\xi_2;h) a(x, \xi) \overline{a(x, \eta)} dy_1 d\xi dw_1 d\xi_1 d\xi_2 dy_2 d\eta dt ds + O(h^\infty).
\end{align}

One can apply stationary phase in $(y_1, \xi)$-variables in \eqref{main kernel1}. The critical point equations for $(y_1, \xi)$ are
\begin{align}
\xi&=\xi_1+ t\partial_{y_1} p_1(y_1,\xi_1)+ O(t^2), \\
y_1&=x.
\end{align}

With applying Taylor expansion at the critical point $(y_{1c}, \xi_{c})$ one has that \eqref{main kernel1} equals
\begin{align}
I (x,x; h):=&h^{-3n}\int \hat{\rho}(s) \hat{\rho}(t)  \exp\Big[\frac{i}h\Big(\left<x-w_1, \xi_1\right>+\left<w_1-y_2,\xi_2\right>+ t(p_1(x,\xi_1)-E_1\big)+\nonumber \\
&s\big(p_2(w_1,\xi_2)-E_2\big) +\left<y_2-x, \eta\right>+O_{w_1,\xi_2}(s^2)+O_{x,\xi_1}(t^2) \Big)\Big]\cdot  \nonumber\\
& c_1(s,t,x, w_1,y_2,\xi_1,\xi_2;h)[a(x, \xi_1)+o(1)] \overline{a(x, \eta)} dw_1 d\xi_1 d\xi_2 dy_2 d\eta dt ds,   \label{main kernel2}
\end{align}
here $c_1\sim \sum_{j=0}^\infty c_{1,j} h^j$ and $c_{1,j}\in C^\infty$.

One can apply stationary phase in $(y_2, \eta)$-variables in \eqref{main kernel2}. The critical point equations for $(y_2, \eta)$ are
\begin{align}
\eta&=\xi_2, \\
y_2&=x.
\end{align}
With applying stationary phase at the critical point $(y_{2c}, \eta_{c})$ and Taylor expansion one has that \eqref{main kernel2} equals
\begin{align*}
I (x,x; h)=&h^{-2n}\int \hat{\rho}(s) \hat{\rho}(t)  \exp\Big[\frac{i}h\Big(\left<x-w_1, \xi_1-\xi_2\right>+ t(p_1(x,\xi_1)-E_1\big)+s\big(p_2(w_1,\xi_2)-E_2\big)  \nonumber \\
+&O_{w_1,\xi_2}(s^2)+O_{x,\xi_1}(t^2) \Big)\Big]  c_2(s,t,x, w_1,\xi_1,\xi_2;h) [a(x, \xi_1)+o(1)] \overline{a(x, \xi_2)}  dw_1 d\xi_1 d\xi_2 dt ds,
\end{align*}
here $c_2\sim \sum_{j=0}^\infty c_{2,j} h^j$ and $c_{2,j}\in C^\infty$.

Next one would like to apply stationary phase in $(w_1, \xi_2)$-variables. The critical point equations for $(w_1, \xi_2)$ are
\begin{align}
\xi_2&=\xi_1-s\partial_{w_1} p_2(w_1,\xi_2)+ O(s^2), \\
w_1&=x-s\partial_{\xi_2} p_2(w_1,\xi_2)+ O(s^2).
\end{align}

Hence with applying Taylor expansion one can get that
\begin{align*}
I (x,x; h)=&h^{-n}\int \hat{\rho}(s) \hat{\rho}(t)  \exp\Big[\frac{i}h\Big(  t(p_1(x,\xi_1)-E_1\big)+s\big(p_2(x,\xi_1)-E_2\big)+O_{x,\xi_1}(t^2)  \nonumber \\
&+O_{x,\xi_1}(s^2) \Big)\Big]  c_3(s,t,x,\xi_1;h) [a(x, \xi_1)+o(1)] [\overline{a(x, \xi_1)}+o(1)]  d\xi_1 dt ds.
\end{align*}
here $c_3\sim \sum_{j=0}^\infty c_{3,j} h^j$ and $c_{3,j}\in C^\infty$.


Notice that $P_1(h)=-h^2\Delta_g$ with $E_1=1$. We use the geodesic normal coordinate about $x$ and can make the change of variables $\xi_1=r\Theta$, where $r=|\xi_1|$ and $\Theta=(\theta_1,\theta_2,\dots,\theta_{n-1})\in S^{n-1}$. One can get
\begin{align}\label{tilde I}
I &(x, x; h)= h^{-n}\int\int\int\int \hat{\rho}(s) \hat{\rho}(t)  \exp\Big[\frac{i}h\Big( t (r^2-1\big)+s\big(p_2(x, r\Theta)-E_2\big)+O_{x,\xi_1}(t^2) \nonumber\\ 
&+O_{x,\xi_1}(s^2) \Big)\Big]c_3(s,t, x,r\Theta;h)[a(x, r\Theta)+o(1)] [\overline{a(x, r\Theta)}+o(1)] r^{n-1} dr d\Theta dt ds.
\end{align}

Again one would like to apply stationary phase in $(r, t)$-variables. The critical point $(r_c,t_c)$ satisfies that
\begin{align}
2r t  +s\Theta\cdot \partial_{\xi_1} p_2(x, r\Theta)+O_{x,\xi_1}(t^2)+O_{x,\xi_1}(s^2)&=0, \label{r}\\
(r^2-1)+O_{x,\xi_1}(t)&=0. \label{E_1}
\end{align}

Taking $\epsilon_0$ small enough in the definition of $\rho$ and combining \eqref{r} and \eqref{E_1}, we can get that
\begin{equation}\label{t_c}
t_c=O(s).
\end{equation}


Taking derivative about $\Theta$ in \eqref{r}, \eqref{E_1} and combining \eqref{E_1} one can derive that
\begin{equation}\label{d t_c}
\partial_{\theta_j} t_c =O(s).
\end{equation}

Now performing stationary phase at the critical point $\big(r_c, t_c \big)$ in \eqref{tilde I} and Taylor expansion gives,
\begin{align}
I (x, x; h)=&h^{-n+1}\int\int \hat{\rho}(s)  \exp\Big[\frac{i}h s\big(p_2(x, \Theta)-E_2\big)+O_{x,\Theta}(st_c) +O_{x,\Theta}(t_c^2)+O_{x,\Theta}(s^2) \Big] \nonumber\\ 
& \cdot c_4(s, x, \Theta; h) [a(x, \Theta)+o(1)] [\overline{a(x, \Theta)}+o(1)] d\Theta ds,
\end{align}
here $c_4\sim \sum_{j=0}^\infty c_{4,j} h^j$ and $c_{4,j}\in C^\infty$.


{\bf Case I}. %\cs Finally, using the Morse condition and separate the range of $s$, one has that
%\begin{equation}
%I (x, x; h)=O(1) |a(x,\Theta_0)|^2 I_h(x)^2,
%\end{equation}
%where $\partial_{\Theta} p_2(x,\Theta_0)=0$.
If $|s|\leq h$, it's straightforward to get that
\begin{equation}
\Big| I (x, x; h) \Big|= O(h^{-n+2}).
\end{equation}

If $|s|>h$, one can perform the stationary phase theorem about $\Theta$ due to the assumption that $(p_1, p_2)$ is Morse type. Hence
\begin{equation}
\Big| I (x, x; h) \Big|\leq C |a(x,\xi_0)|^2 h^{\frac{-n+1}2} \Big|\int_{|s|>h}s^{\frac{-n-1}2} ds\Big|=O(1) |a(x,\xi_0)|^2 I_h(x)^2.
\end{equation}


{\bf Case II}. If $(p_1|_{\Gamma}, p_2|_{\Gamma})$ is {\bf rank $1$} at $x$, integration by parts in some variable $\theta_j$ is necessary to deal with integration $|{I} (x,x; h)|$. 

Now one can separate the integration
\begin{align*}
\Big| I (x, x; h) \Big|& \leq \Big|h^{-n+1}\int_{|s|\leq h} \int \hat{\rho}(s)  \exp[\frac{i}h\Phi(s,x,\Theta) ] \cdot c_5(s, x, \Theta; h)d\Theta ds \Big| \nonumber \\
+& \Big|h^{-n+1}\int_{h<|s|<1}\int \hat{\rho}(s)  \exp[\frac{i}h\Phi(s,x,\Theta) ] \cdot c_5(s, x, \Theta; h)d\Theta ds\Big|
:=\Big| {I}_1 (x; h)\Big| + \Big| {I}_2 (x; h) \Big|,
\end{align*}
here we set $\Phi(s,x,\Theta)= s\big(p_2(x, \Theta)-E_2\big)+O_{x,\Theta}(st_c) +O_{x,\Theta}(t_c^2)+O_{x,\Theta}(s^2)$ and $c_5(s, x, \Theta; h)=c_4(s, x, \Theta; h) [a(x, \Theta)+o(1)] [\overline{a(x, \Theta)}+o(1)]$.

\noindent It's straightforward to get that 
\begin{equation}\label{last1}
\Big| {I}_1 (x; h)\Big|=O(h^{-n+2}).
\end{equation}

\noindent %One sets an operator
%\begin{equation*}
%L:=\frac{1}{|\nabla_\Theta \Phi|^2}\sum_{j=1}^{n-1} \partial_{\Theta_j}\Phi\,h\, D_{\Theta_j}.
%\end{equation*}
For the second term, with the help of \eqref{t_c} and \eqref{d t_c} the key observation is that for some $1\leq j\leq n-1$
\begin{equation}\label{phi low}
\big| \partial_{\theta_j} \Phi(s,x,\Theta) \big|\geq|s| \big(\big|\partial_{\theta_j} p_2(x, \Theta)\big|-O(\epsilon_0)\big)>C|s|>0
\end{equation}
due to the {\bf rank $1$ condition}  and the definition of $\rho$.

Also with the help of \eqref{t_c} and \eqref{d t_c} one has that
\begin{equation}\label{phi up}
\big|\partial_{\theta_j}^\alpha \Phi(s,x,\Theta)\big|\leq C|s|, \quad\text{for a given multi-index } \alpha.
\end{equation}


Hence, one can integrate by parts to get that
\begin{align}\label{last2}
\Big| {I}_2 (x; h) \Big|&=h^{-n+2} \Big|\int_{h<|s|<1} \int \frac{1}{\partial_{\theta_j} \Phi} \big(\partial_{\theta_j} \exp[\frac{i}h  \Phi(s,x,\Theta)]\big)  \cdot c_5(s, x, \Theta; h) d\Theta ds\Big| \nonumber\\
&\leq h^{-n+2}\Big|\int_{h<|s|<1} \int \frac{1}{\partial_{\theta_j} \Phi} \exp[\frac{i}h  \Phi(s,x,\Theta)]  \cdot \partial_{\theta_j} c_5(s, x, \Theta; h) d\Theta ds\Big| \nonumber\\
&+ h^{-n+2} \Big|\int_{h<|s|<1} \int \frac{\partial^2_{\theta_j}\Phi }{|\partial_{\theta_j} \Phi|^2}   \exp[\frac{i}h  \Phi(s,x,\Theta)] c_5(s, x, \Theta; h) d\Theta ds\Big|. \nonumber
\end{align}

One more integration by parts gives us that
\begin{align}
\Big| {I}_2 (x; h) \Big|
&\leq h^{-n+3}\Big|\int_{h<|s|<1} \int \frac{1}{|\partial_{\theta_j} \Phi|^2} \big(\partial_{\theta_j} \exp[\frac{i}h  \Phi(s,x,\Theta)]\big)  \cdot \partial_{\theta_j}c_5(s, x, \Theta; h) d\Theta ds\Big| \nonumber\\
&+ h^{-n+3} \Big|\int_{h<|s|<1} \int \frac{\partial^2_{\theta_j}\Phi }{(\partial_{\theta_j} \Phi)^3}  \big(\partial_{\theta_j} \exp[\frac{i}h  \Phi(s,x,\Theta)]\big) c_5(s, x, \Theta; h) d\Theta ds\Big| \nonumber\\
&\leq h^{-n+3}\Big|\int_{h<|s|<1} \int \frac{1}{|\partial_{\theta_j} \Phi|^2} \exp[\frac{i}h  \Phi(s, x, \Theta; h)]  \cdot \partial^2_{\theta_j}c_5(s, x, \Theta; h) \, d\Theta ds\Big| \nonumber\\
&+2 h^{-n+3}\Big|\int_{h<|s|<1} \int \frac{\partial_{\theta_j}^2\Phi}{(\partial_{\theta_j} \Phi)^3}  \exp[\frac{i}h  \Phi(s, x, \Theta; h)] c_5(s, x, \Theta; h) d\Theta ds\Big| \nonumber\\
&+h^{-n+3}\Big|\int_{h<|s|<1} \int \frac{\partial_{\theta_j}^2\Phi}{(\partial_{\theta_j} \Phi)^3}  \exp[\frac{i}h  \Phi(s, x, \Theta; h)] \cdot \partial_{\theta_j}c_5(s, x, \Theta; h) d\Theta ds\Big| \nonumber\\
&+ h^{-n+3} \Big|\int_{h<|s|<1} \int \big( \frac{\partial^3_{\theta_j}\Phi }{(\partial_{\theta_j} \Phi)^3}-3 \frac{(\partial^2_{\theta_j}\Phi)^2 }{(\partial_{\theta_j} \Phi)^4} \big)  \exp[\frac{i}h  \Phi(s,x,\Theta)] c_5(s, x, \Theta; h) d\Theta ds\Big| \nonumber\\
&\leq C h^{-n+3}  \Big|\int_{h<|s|<1}\frac{1}{s^2} ds \Big|=Ch^{-n+2},
\end{align}
with the help of \eqref{phi low} and \eqref{phi up} in the second last inequality.
\end{proof}

\bigskip
Now let us take a real-valued function $\chi\in \mathcal{S}(\mathbb{R})$ satisfying
\begin{equation}\label{spt}
\chi(0)=1\quad \text{and}\quad \hat\chi(t)=0,\quad |t|\geq\epsilon,
\end{equation}
here $\epsilon$ is a small positive constant.
Since $\chi\big(h^{-1}[P_1(h)-E_1(h)]\big)u_h=u_h$, by the Cauchy-Schwarz inequality, one has that
\begin{align}\label{red1}
&\Big| A u_h (x) \Big|^2 \nonumber \\
=&\Big| \Big(A \chi \big( h^{-1}[P_1(h)-E_1(h)]\big) u_h\Big) (x) \Big|^2 \nonumber\\
=&  \Big| \int \Big(A\chi\big( h^{-1}[P_1(h)-E_1(h)]\big)\Big) (x,y) u_h(y) dy  \Big|^2  \nonumber \\
\leq &  \int \Big| A\chi\big(h^{-1}[P_1(h)-E_1(h)]\big) (x,y) \Big|^2 dy.
\end{align}

Furthermore, using the orthogonality of $\{u_j^h\}$, we can deduce that
\begin{align}\label{red2}
&\int \Big| A\chi\big(h^{-1}[P_1(h)-E_1(h)]\big) (x,y) \Big|^2 dy \nonumber\\
=& \int \sum_j \chi(h^{-1}[\lambda_j^{(1)}(h)-E_1(h)]) A u_j^h(x) \overline{u_j^h(y)} \cdot \sum_l \chi(h^{-1}[\lambda_l^{(1)}(h)-E_1(h)]) \overline{A u_l^h(x)} u_l^h(y) dy \nonumber \\
=& \sum_j \rho(h^{-1}[\lambda_j^{(1)}(h)-E_1(h)]) A u_j^h(x) \overline{A u_j^h(x)}.
\end{align}
with setting $\rho(t)=(\chi(t))^2$. 

We claim that
\begin{align}\label{upshot}
&\sum_j \rho(h^{-1}[\lambda_j^{(1)}(h)-E_1(h)]) A u_j^h(x) \overline{A u_j^h(x)} \leq 
O(1)\cdot \nonumber \\ &\sup_{\{E_2;\, (E_1,E_2)\in \mathcal{P}(T^*M)\}}\sum_j \rho(h^{-1}[\lambda_j^{(1)}(h)-E_1]) 
\rho(h^{-1}[\lambda_j^{(2)}(h)-E_2])  A u^h_j (x) \overline{A u^h_j(x)}.
\end{align}


%With noticing that $\rho \big(h^{-1}[P_2(h)-E_2(h)]\big) u_j^h=u_j^h$, one can get that
%\begin{align}\label{main kernel}
%&\sum_j \rho(h^{-1}[\lambda_j^{(1)}(h)-E_1(h)]) A u_j^h(x) \overline{A u_j^h(x)} \nonumber\\
%=& \sum_j \rho(h^{-1}[\lambda_j^{(1)}(h)-E_1(h)]) A \big(\rho \big(h^{-1}[P_2(h)-E_2(h)]\big) u_j^h\big)(x) %\overline{A u_j^h(x)} \nonumber\\
%=& \sum_j \rho(h^{-1}[\lambda_j^{(1)}(h)-E_1(h)]) A \Big(\rho \big(h^{-1}[\lambda_j^{(2)}(h)-E_2(h)]\big) u_j^h(x) \Big) \overline{A u_j^h(x)} \nonumber\\
%=& \sum_j \rho(h^{-1}[\lambda_j^{(1)}(h)-E_1(h)]) \rho \big(h^{-1}[\lambda_j^{(2)}(h)-E_2(h)]\big) A u_j^h(x) \overline{A u_j^h(x)} \nonumber\\
%=&  \frac{1}{2\pi} \int \hat\rho(t) \Big(A e^{\frac{i}h t(P_1-E_1)} \rho\big(h^{-1}[P_2(h)-E_2(h)]\big)A^* \Big)(t,x,x) dt.
%\end{align}
%In conclusion, one has that
%\begin{equation}
%\Big| A u_h (x) \Big|^2\leq \frac{1}{2\pi} \int \hat\rho(t) \Big(A e^{\frac{i}h t(P_1-E_1)} \rho\big(h^{-1}[P_2(h)-E_2(h)]\big)A^* \Big)(t,x,x) dt
%\end{equation}



%Hence proving \eqref{main eq} is equivalent to showing that
%\begin{equation}\label{upshot}
%\Big|  \sum_j \rho(h^{-1}[\lambda_j^{(1)}(h)-E_1(h)]) u^h_j(x) \overline{u^h_j(x)} \Big| \leq C h^{-n+2} |\ln h| .
%\end{equation}


%We claim that we need to show that
%\begin{align}\label{upshot5}
%\sup_{\{E_2;\, (E_1,E_2)\in \mathcal{P}(T^*M)\}}\sum_j \rho(h^{-1}[\lambda_j^{(1)}(h)-E_1])\rho(h^{-1}[\lambda_j^{(2)}(h)-E_2])  u^h_j (x) \overline{u^h_j(x)} 
%\leq C h^{-n+2} |\ln h|.
%\end{align}

Indeed from \cite[Section 2]{Tot} one knows that there exists a constant $C_2>0$ (independent of $j$ and $h$) such that for any $h\in(0, h_0]$, and $(\lambda_j^{(1)}(h), \lambda_j^{(2)}(h))\in \text{Spec}(P_1(h), P_2(h))$, with $|\lambda_j^{(1)}(h)-E_1|\leq C_1 h$,
\begin{equation*}
\inf_{E_2}|\lambda_j^{(2)}(h)-E_2|\leq C_2 h.
\end{equation*}
So, after taking $\epsilon$ in \eqref{spt} sufficiently small there exists a constant $C_3>0$ (independent of $j$ and $h$) such that for all $j\geq 1$ and $h\in(0, h_0]$,
\begin{equation*}
\sup_{\{E_2;\, (E_1,E_2)\in \mathcal{P}(T^*M),\, |\lambda_j^{(1)}(h)-E_1|\leq C_1 h\}}\rho\big(h^{-1}(\lambda_j^{(2)}(h) -E_2)\big)\geq C_3>0.
\end{equation*}
Since the sum on the RHS of \eqref{upshot} has non-negative terms, by restricting to $\{j;\,|\lambda_j^{(1)}(h)-E_1|\leq C_1 \}$ and (after taking $\epsilon$ small enough) using that $\rho(h^{-1}[\lambda_j^{(1)}(h)-E_1])\geq C_4>0$ for these eigenvalues, one gets that
\begin{align*}
\sum_{\{j;\, |\lambda_j^{(1)}(h)-E_1|=O(h)\}} A u^h_j(x) \overline{Au^h_j(x)} \leq 
O(1)\sup_{\{E_2;\, (E_1,E_2)\in \mathcal{P}(T^*M)\}}& \sum_j \rho(h^{-1}[\lambda_j^{(1)}(h)-E_1])\cdot \nonumber \\  
&\rho(h^{-1}[\lambda_j^{(2)}(h)-E_2])  A u^h_j (x) \overline{A u^h_j(x)},
\end{align*}
which can deduce \eqref{upshot}.

With noticing that
\begin{align}\label{sum to integral}
&\sum_j \rho(h^{-1}[\lambda_j^{(1)}(h)-E_1]) \rho \big(h^{-1}[\lambda_j^{(2)}(h)-E_2]\big) A u_j^h(x) \overline{A u_j^h(x)} \nonumber\\
=&  \frac{1}{2\pi} \int \hat\rho(t) \Big(A e^{\frac{i}h t(P_1-E_1)} \rho\big(h^{-1}[P_2(h)-E_2]\big)A^* \Big)(t,x,x) dt,    
\end{align}
and combining \eqref{red1}, \eqref{red2} and \eqref{upshot}, one can apply Lemma \ref{lwl} to get following,

\begin{cor}\label{cor}
Under the same assumption of Lemma \ref{lwl}, if $(p_1|_{\Gamma}, p_2|_{\Gamma})$ with $\Gamma:=\mathcal{C}_x\cap \supp a$ is {\bf rank $1$} at $x$, one has that 
\begin{equation}\label{main cor eq2}
|A(x,hD) u_{h}|^2=O(h^{-n+2}).
\end{equation}
Otherwise, one has the following
\begin{equation}
|A(x,hD) u_{h}|^2=O(1) I_h(x)^2,
\end{equation}
if $(p_1, p_2)$ is Morse type.
\end{cor}

Furthermore, we can prove that
\begin{lem}\label{lwl 2}
Assume $A(x, hD)\in \Psi^0(M)$ with symbol $a=\sum_{j=0}^\infty a_j h^j$. Let $\rho\in C_0^\infty(\mathbb{R})$ vanish for $|t|\geq \epsilon_0$, $\epsilon_0>0$ small enough, and satisfy $\rho(0)=1$. Let $\Gamma:=\mathcal{C}_x\cap \supp a$. Then if $(p_1|_\Gamma, p_2|_\Gamma, p_3|_\Gamma)$ is {\bf rank $2$} at $x$, one has that
\begin{equation}\label{main lemma 2 eq}
\int_{\mathbb{R}} \hat{\rho}(t) \big(A e^{\frac{i}h t(P_1-E_1)}  \rho\big(h^{-1}[P_2(h)-E_2]\big)\rho\big(h^{-1}[P_3(h)-E_3]\big)  A^*\big)(t,x,x) dt =O(h^{-n+3}).
\end{equation}
\end{lem}
\begin{proof}
With repeating the begining of the proof of Lemma \ref{lwl}, one need to show that the right hand of \eqref{main lemma 2 eq} can bound
\begin{align}
I(x,x;h):=h^{- 6n}\int &\hat{\rho}(s_2)\hat{\rho}(s_1)\hat{\rho}(t) \exp[\frac{i}h \left<x-y_1, \xi\right>] \exp[\frac{i}h \Phi(s_1, s_2,t,y_1, w_1, w_2, y_2, \xi_1, \xi_2, \xi_3)]\cdot  \nonumber \\
& \exp[\frac{i}h \left<y_2-x, \eta\right>] c_0(s_1, s_2,t,y_1,w_1,w_2,y_2,\xi_1,\xi_2,\xi_3;h) a(x, \xi) \cdot \nonumber \\
 &\qquad\qquad \overline{a(x, \eta)} dy_1 d\xi dw_1 d\xi_1 dw_2 d\xi_2 d\xi_3 dy_2 d\eta dt ds_1 ds_2,
\end{align}
here $c_0\sim \sum_{j=0}^\infty c_{0,j} h^j$ with $c_{0,j}\in C^\infty$ and $c$ is supported near the energy shell $\{p_1=E_1\}\cup\{p_2=E_2\}\cup\{p_3=E_3\}$, 
\begin{align*}
\Phi(\cdot)=&\left<y_1-w_1,\xi_1 \right>+\left<w_1-w_2, \xi_2\right>+\left<w_2-y_2, \xi_3\right> +t\big(p_1(y_1,\xi_1)-E_1\big)+s_1\big( p_2(w_1,\xi_2)-E_2 \big)\nonumber \\
&+s_2\big( p_3(w_2,\xi_3)-E_3 \big)+O_{y_1,\xi_1}(t^2)+O_{w_1,\xi_2}(s_1^2)+O_{w_2,\xi_3}(s_2^2).
\end{align*}

Now apply stationary phase in $(y_1, \xi)$-variables. The critical point equations for $(y_1, \xi)$ are
\begin{align}
\xi&=\xi_1+ t\partial_{y_1} p_1(y_1,\xi_1)+ O(t^2), \\
y_1&=x.
\end{align}

With applying Taylor expansion at the critical point $(y_{1c}, \xi_{c})$ one has that
\begin{align}
&I (x,x; h):=h^{-4n}\int \hat{\rho}(s_2)\hat{\rho}(s_1)\hat{\rho}(t)  \exp\Big[\frac{i}h\Big(\left<x-w_1, \xi_1\right>+\left<w_1-w_2,\xi_2\right>+\left<w_2-y_2,\xi_3\right> \nonumber \\
&\qquad+\left<y_2-x, \eta\right> +t(p_1(x,\xi_1)-E_1\big)+ s_1\big(p_2(w_1,\xi_2)-E_2\big) +s_2\big( p_3(w_2,\xi_3)-E_3 \big)+O_{x,\xi_1}(t^2) \nonumber\\
&\qquad\qquad + O_{w_1,\xi_2}(s_1^2)+O_{w_2,\xi_3}(s_2^2) \Big)\Big]\cdot c_1(s_1, s_2,t,x, w_1,w_2,y_2,\xi_1,\xi_2,\xi_3;h)[a(x, \xi_1)+o(1)]  \nonumber\\
&\qquad\qquad\qquad \cdot \overline{a(x, \eta)} dw_1 d\xi_1 dw_2 d\xi_2 d\xi_3 dy_2 d\eta dt ds_1 ds_2, 
\end{align}
here $c_1\sim \sum_{j=0}^\infty c_{1,j} h^j$ and $c_{1,j}\in C^\infty$.

Apply stationary phase in $(y_2, \eta)$-variables. The critical point equations for $(y_2, \eta)$ are
\begin{align}
\eta&=\xi_3, \\
y_2&=x.
\end{align}
With applying stationary phase at the critical point $(y_{2c}, \eta_{c})$ and Taylor expansion one has that
\begin{align*}
I (x,x; h)=&h^{-3n}\int \hat{\rho}(s_2)\hat{\rho}(s_1)\hat{\rho}(t)  \exp\Big[\frac{i}h\Big(\left<x-w_1, \xi_1\right>+\left<w_1-w_2,\xi_2\right>+\left<w_2-x,\xi_3\right> \nonumber \\
&\quad +t(p_1(x,\xi_1)-E_1\big)+ s_1\big(p_2(w_1,\xi_2)-E_2\big) +s_2\big( p_3(w_2,\xi_3)-E_3 \big)+O_{x,\xi_1}(t^2) \nonumber\\
&\quad\quad + O_{w_1,\xi_2}(s_1^2)+O_{w_2,\xi_3}(s_2^2) \Big)\Big]\cdot c_2(s_1, s_2,t,x, w_1,w_2,\xi_1,\xi_2,\xi_3;h) [a(x, \xi_1)+o(1)]  \nonumber\\
&\quad\qquad\quad \cdot \overline{a(x, \xi_3)} dw_1 d\xi_1 dw_2 d\xi_2 d\xi_3 dt ds_1 ds_2, 
\end{align*}
here $c_2\sim \sum_{j=0}^\infty c_{2,j} h^j$ and $c_{2,j}\in C^\infty$.

Apply stationary phase in $(w_1, \xi_2)$-variables. The critical point equations for $(w_1,\xi_2)$ are
\begin{align}
\xi_2&=\xi_1-s_1\partial_{w_1} p_2(w_1,\xi_2)+ O_{w_1,\xi_2}(s_1^2), \\
%\xi_3&=\xi_2-s_2\partial_{w_2} p_3(w_2,\xi_3)+ O_{w_2,\xi_3}(s_2^2),\\
w_1&=w_2-s_1\partial_{\xi_2} p_2(w_1,\xi_2)+ O_{w_1, \xi_2}(s_1^2).
%w_2&=x-s_2\partial_{\xi_3} p_3(w_2,\xi_3)+ O_{w_2,\xi_3}(s_2^2).
\end{align}

Hence with applying Taylor expansion one can get that
\begin{align*}
I (x,x; h)=&h^{-2n}\int \hat{\rho}(s_2)\hat{\rho}(s_1)\hat{\rho}(t)  \exp\Big[\frac{i}h\Big(\left<x-w_2, \xi_1-\xi_3\right>+t(p_1(x,\xi_1)-E_1\big)+ \nonumber \\
&  s_1\big(p_2(w_2,\xi_1)-E_2\big)+s_2\big( p_3(w_2,\xi_3)-E_3 \big)+O_{x,\xi_1}(t^2)+ O_{w_2,\xi_1}(s_1^2)+O_{w_2,\xi_3}(s_2^2) \Big)\Big] \nonumber\\
&\qquad \cdot c_3(s_1,s_2,t,x,w_2,\xi_1, \xi_3;h)[a(x, \xi_1)+o(1)] \overline{a(x, \xi_3)} d\xi_1 dw_2 d\xi_3 dt ds_1 ds_2,
\end{align*}
here $c_3\sim \sum_{j=0}^\infty c_{3,j} h^j$ and $c_{3,j}\in C^\infty$.

Next apply stationary phase in $(w_2, \xi_3)$-variables. The critical point equations for $(w_2,\xi_3)$ are
\begin{align}
\xi_3&=\xi_1-s_1\partial_{w_2} p_2(w_2,\xi_1)-s_2\partial_{w_2} p_3(w_2,\xi_3)+O_{w_2,\xi_1}(s_1^2)+ O_{w_2,\xi_3}(s_2^2),\\
w_2&=x-s_2\partial_{\xi_3} p_3(w_2,\xi_3)+ O_{w_2,\xi_3}(s_2^2).
\end{align}

So applying Taylor expansion one can get that
\begin{align}
I (x,x; h)=&h^{-n}\int \hat{\rho}(s_2)\hat{\rho}(s_1)\hat{\rho}(t)  \exp\Big[\frac{i}h\Big(t(p_1(x,\xi_1)-E_1\big)+ s_1\big(p_2(x,\xi_1)-E_2\big)+ \nonumber \\
&\qquad  s_2\big( p_3(x,\xi_1)-E_3 \big)+O_{x,\xi_1}(t^2)+O_{x,\xi_1}(s_1 s_2)+ O_{x,\xi_1}(s_1^2)+O_{x,\xi_1}(s_2^2) \Big)\Big] \nonumber\\
&\qquad\qquad \cdot c_4(s_1,s_2,t,x,\xi_1;h) [a(x, \xi_1)+o(1)] [\overline{a(x, \xi_1)}+o(1)] d\xi_1 dt ds_1 ds_2,
\end{align}
here $c_4\sim \sum_{j=0}^\infty c_{4,j} h^j$ and $c_{4,j}\in C^\infty$.

Now use the geodesic normal coordinate about $x$ and make the change of variables $\xi_1=r\Theta$, where $r=|\xi_1|$ and $\Theta=(\theta_1,\theta_2,\dots,\theta_{n-1})\in S^{n-1}$. One can get
\begin{align}\label{tilde I'}
I (x, x; h)=& h^{-n}\int \hat{\rho}(s_2)\hat{\rho}(s_1)\hat{\rho}(t)  \exp\Big[\frac{i}h\Big( t (r^2-1\big)+s_1\big(p_2(x, r\Theta)-E_2\big)+s_2\big( p_3(x,r\Theta)-E_3 \big) \nonumber\\ 
&+O_{x,\xi_1}(t^2)+O_{x,\xi_1}(s_1 s_2)+ O_{x,\xi_1}(s_1^2)+O_{x,\xi_1}(s_2^2) \Big)\Big]c_4(s_1,s_2,t, x,r\Theta;h)\nonumber \\
&\quad\cdot[a(x, r\Theta)+o(1)] [\overline{a(x, r\Theta)}+o(1)] r^{n-1} dr d\Theta dt ds_1 ds_2.
\end{align}

Again one would like to apply stationary phase in $(r, t)$-variables. The critical point $(r_c,t_c)$ satisfies that
\begin{align}
2r t  +s_1\Theta\cdot \partial_{\xi_1} p_2(x, r\Theta)+s_2\Theta\cdot \partial_{\xi_1} p_3(x, r\Theta)+O_{x,\xi_1}(t^2)+O_{x,\xi_1}(s_1 s_2)\qquad\qquad & \nonumber \\
+ O_{x,\xi_1}(s_1^2)+O_{x,\xi_1}(s_2^2)&=0, \label{2r}\\
(r^2-1)+O_{x,\xi_1}(t)&=0. \label{2E_1}
\end{align}

Taking $\epsilon_0$ small enough in the definition of $\rho$ and combining \eqref{2r} and \eqref{2E_1}, we can get that
\begin{equation}\label{t_c 2}
t_c=O(s_1)+O(s_2).
\end{equation}


Taking derivative about $\theta_j$ in \eqref{2r}, \eqref{2E_1} and combining \eqref{2E_1} one can derive that
\begin{equation}\label{d t_c 2}
\partial_{\theta_j} t_c =O(s_1)+O(s_2).
\end{equation}

Now performing stationary phase at the critical point $\big(r_c, t_c \big)$ in \eqref{tilde I'} and Taylor expansion gives,
\begin{align}
I (x, x; h)=&h^{-n+1}\int \hat{\rho}(s_1)\hat{\rho}(s_2)  \exp\Big[\frac{i}h \Big(s_1\big(p_2(x, \Theta)-E_2\big)+s_2\big(p_3(x, \Theta)-E_3\big)+O_{x,\Theta}(s_1t_c) \nonumber\\ 
& +O_{x,\Theta}(s_2t_c)+O_{x,\Theta}(t_c^2)+O_{x,\Theta}(s_1 s_2)+O_{x,\Theta}(s_1^2)+O_{x,\Theta}(s_2^2)\Big) \Big] c_5(s_1, s_2, x, \Theta; h)\nonumber\\
&\cdot [a(x, \Theta)+o(1)] [\overline{a(x, \Theta)}+o(1)] d\Theta ds_1 ds_2,
\end{align}
here $c_5\sim \sum_{j=0}^\infty c_{5,j} h^j$ and $c_{5,j}\in C^\infty$.

Then one can separate the integration
\begin{align*}
\Big| I (x, x; h) \Big|& \leq \Big|h^{-n+1}\int_{|s_1|,|s_2|\leq h} \int \hat{\rho}(s_1)\hat{\rho}(s_2)  \exp[\frac{i}h\Phi(s_1,s_2,z,\Theta) ] \cdot c_6(s_1, s_2, x, \Theta; h)d\Theta ds_1 ds_2 \Big| \nonumber \\
+&\Big|h^{-n+1}\int_{|s_1|\leq h,h<|s_2|<1}\int \cdot\, d\Theta ds_1ds_2\Big|+\Big|h^{-n+1}\int_{h<|s_1|<1,|s_2|\leq h}\int \cdot\, d\Theta ds_1ds_2\Big| \nonumber\\
+& \Big|h^{-n+1}\int_{h<|s_1|,|s_2|<1}\int \hat{\rho}(s_1)\hat{\rho}(s_2)  \exp[\frac{i}h\Phi(s_1,s_2,z,\Theta) ] \cdot c_6(s_1,s_2, x, \Theta; h)d\Theta ds_1ds_2\Big| \nonumber\\
&:=\Big| {I}_1 (x; h)\Big| + \Big| {I}_2 (x; h) \Big|+ \Big| {I}_3 (x; h) \Big|+ \Big| {I}_4 (x; h) \Big|,
\end{align*}
here we set $\Phi(s_1,s_2,x,\Theta)= s_1\big(p_2(x, \Theta)-E_2\big)+s_2\big(p_3(x, \Theta)-E_3\big)+O_{x,\Theta}(s_1t_c) +O_{x,\Theta}(s_2t_c)+O_{x,\Theta}(t_c^2)+O_{x,\Theta}(s_1 s_2)+O_{x,\Theta}(s_1^2)+O_{x,\Theta}(s_2^2)$ and $c_6(s_1,s_2, x, \Theta; h)=c_5(s_1, s_2, x, \Theta; h) [a(x, \Theta)+o(1)] [\overline{a(x, \Theta)}+o(1)]$.

We only focus on $\Big| {I}_1 (x; h) \Big|$ and $\Big| {I}_4 (x; h) \Big|$ since the rest can be dealt with similarly. It's straightforward to get that 
\begin{equation}\label{last1'}
\Big| {I}_1 (z; h)\Big|=O(h^{-n+3}).
\end{equation}


For $\Big| {I}_4 (x; h) \Big|$, %one sets
%\begin{equation*}
%\Phi(s_1,s_2,x,\Theta)=\Phi_1(s_1,s_2,x,\Theta)+\Phi_2(s_1,s_2,x,\Theta),
%\end{equation*}
%where
%\begin{align*}
%\Phi_1(s_1,s_2,x,\Theta)&=s_1\big(p_2(x, \Theta)-E_2\big)+O_{x,\Theta}(s_1t_c) +O_{x,\Theta}(t_c^2)+O_{x,\Theta}(s_1 s_2)+O_{x,\Theta}(s_1^2)\\
%\Phi_2(s_1,s_2,x,\Theta)&=s_2\big(p_3(x, \Theta)-E_3\big) +O_{x,\Theta}(s_2t_c)+O_{x,\Theta}(t_c^2)+O_{x,\Theta}(s_2^2)
%\end{align*}
with the help of \eqref{t_c 2} and \eqref{d t_c 2} the key observation is that for some $1\leq j\neq k\leq n-1$
\begin{align}
\big| \partial_{\theta_j} \Phi(s_1,s_2,x,\Theta) \big|&\geq|s_1| \big(\big|\partial_{\theta_j} p_2(x, \Theta)\big|-O(\epsilon_0)\big)>C_1|s_1|>0, \label{rank2 1}\\
\big| \partial_{\theta_k} \Phi(s_1,s_2,x,\Theta) \big|&\geq|s_2| \big(\big|\partial_{\theta_k} p_3(x, \Theta)\big|-O(\epsilon_0)\big)>C_2|s_2|>0 \label{rank2 2}
\end{align}
due to the {\bf rank $2$ condition}  and the definition of $\rho$.

Also with the help of \eqref{t_c 2} and \eqref{d t_c 2} one has that
\begin{equation}\label{s1s2}
\big|\partial_{\Theta}^\alpha \Phi(s_1, s_2, x,\Theta)\big|=O(s_1)+O(s_2), \quad\text{for a given multi-index } \alpha.
\end{equation}


Hence, one can adapt the argument from the end of Lemma \ref{lwl}: integrate by parts twice with respect to variables $\theta_j$ and $\theta_k$, then apply equation \eqref{s1s2} to conclude that
\begin{align}\label{last2'}
\Big| {I}_4 (x; h) \Big| \leq C h^{-n+5}  \Big|\int_{h<|s_1|,|s_2|<1}\frac{1}{s^2_1 s^2_2}  ds_1 d_2 \Big|=O(h^{-n+3}) .
\end{align}
%\begin{align}\label{last2}
%\Big| {I}_4 (x; h) \Big|&=h^{-n+2} \Big|\int_{h<|s_1|,|s_2|<1} \int \frac{1}{\partial_{\theta_j} \Phi} \big(\partial_{\theta_j} \exp[\frac{i}h  \Phi(s_1,s_2,x,\Theta)]\big)  \cdot c_6(s_1,s_2, x, \Theta; h) d\Theta ds_1 ds_2\Big| \nonumber\\
%&\leq h^{-n+2}\Big|\int_{h<|s_1|,|s_2|<1} \int \frac{1}{\partial_{\theta_j} \Phi} \exp[\frac{i}h  \Phi(s_1,s_2,x,\Theta)]  \cdot \partial_{\theta_j} c_6(s_1,s_2, x, \Theta; h) d\Theta ds_1 ds_2\Big| \nonumber\\
%&+ h^{-n+2} \Big|\int_{h<|s_1|,|s_2|<1} \int \frac{\partial^2_{\theta_j}\Phi }{|\partial_{\theta_j} \Phi|^2}   \exp[\frac{i}h  \Phi(s_1,s_2,x,\Theta)] c_6(s_1, s_2,x, \Theta; h) d\Theta ds_1 ds_2\Big| \nonumber\\
%&=h^{-n+3}\Big|\int_{h<|s_1|,|s_2|<1} \int \frac{1}{\partial_{\theta_j} \Phi} \exp[\frac{i}h  \Phi(s_1,s_2,x,\Theta)]  \cdot \partial_{\theta_j} c_6(s_1,s_2, x, \Theta; h) d\Theta ds_1 ds_2\Big| \nonumber\\
%&+ h^{-n+3} \Big|\int_{h<|s_1|,|s_2|<1} \int \frac{\partial^2_{\theta_j}\Phi }{|\partial_{\theta_j} \Phi|^2}   \exp[\frac{i}h  \Phi(s_1,s_2,x,\Theta)] c_6(s_1, s_2,x, \Theta; h) d\Theta ds_1 ds_2\Big| \nonumber\\
%&\leq C h^{-n+3}  \Big|\int_{h<|s_1|,|s_2|<1}\frac{1}{s_1 s_2}  ds_1 d_2 \Big|=Ch^{-n+3} |\ln h|^2.
%\end{align}



\end{proof}

Following a similar argument after Lemma \ref{lwl} and combining Lemma \ref{lwl 2}, one can inductively obtain Theorem \ref{trace}.

\bigskip


\section{Proof of Theorem \ref{main1}}\label{section3}
The proof of \eqref{main1 eq2} is an application of Theorem \ref{trace}. Hence one only need to focus on \eqref{main1 eq} and the proof is mainly inspired by \cite{STZ11} with help of the results in Section \ref{section2}. We first introduce some notation. One sets $L_\delta(x,\xi)$, if it exists, is the length of shortest loop such that $\text{dist}(G^{\ell}_x(\xi),\xi)\leq \delta$, with $\text{dist}(\cdot, \cdot)$ being the distance induced by the metric, $\xi$ is the initial direction and $G^{\ell}_x(\xi)$ is the terminal direction of all loop with a given length $\ell$. Let $\mathcal{R}_x^\delta$ be all the $\xi\in \bigcap\limits_{2\leq j \leq n}\{|p_j(x,\xi)-E_j|\leq \epsilon\}$ satisfying $L_\delta(x,\xi)$ is finite.

Without generality, we work on a coordinate patch of a point $x$, which we identified with an open subset of $\mathbb{R}^n$. By continuity, one can choose an open set $U$ such that $\Gamma\subsetneqq U$ and $\mathcal{P}|_{U}$ is rank $k$ ($k\geq 2$ if $n>3$) at $x$, with a reminder that $\Gamma=\mathcal{R}_x\cap \bigcap\limits_{2\leq j \leq n}\{|p_j(x,\xi)-E_j|\leq \epsilon\}$. Then we fix a small number $\delta>0$ so that 
\begin{equation}
\Gamma\subset \mathcal{R}_x^\delta\subset U.
\end{equation}

For given a large number $T$ (to be specified later), applying $C^\infty$ Urysohn lemma, there exists a function $b\in C^\infty(S^{n-1})$ with range $[0,1]$ such that 
\begin{equation}
\mathcal{R}_x^\delta\subset \supp{b}\subset U,
\end{equation}
% \begin{equation}
% \Gamma\cap \supp{b}=\varnothing,
% \end{equation}
and
\begin{equation}
    L_\delta(x,\xi)\geq 2T \quad \text{ for } \forall \xi \in \supp B,
\end{equation}
where $B(\xi)=1-b(\xi)$. 

From Theorem \ref{trace} one has that
\begin{equation}
|b(x, hD)u_h(x)|=o(1)I_n(h).
\end{equation}

So one only need to deal with $|B(x, hD) u_h(x)|$. We construct a smooth partition of unity $1=\sum_k \psi_k(\xi)$ of $S^{n-1}$ which consists of $O(\delta^{-(n-1)})$ terms such that each term has range in $[0,2]$ and is supported in a spherical cap of diameter smaller than $\delta/10$. Let $B_k(x,hD)$ be the zero-order h-pseudo-differential operator whose symbol is $B(x,\xi)\psi_k(\xi)$. Hence
\begin{equation}\label{terms}
    |B(x, hD) u_h(x)| \leq \sum_k \Big| B_k u_h (x) \Big|.
\end{equation}


By Cauchy-Schwarz inequality and using the orthogonality of $\{u_j^h\}$, one has that
\begin{align}
&\Big| B_k u_h (x) \Big|^2 \nonumber \\
=&\Big| \big(B_k \chi \big(T h^{-1}[P_1(h)-E_1]\big) u_h\big) (x) \Big|^2 \nonumber\\
=&  \Big| \int \Big(B_k\chi\big(T h^{-1}[P_1(h)-E_1]\big)\Big) (x,y) u_h(y) dy  \Big|^2  \nonumber \\
\leq &  \int \Big| B_k\chi\big(Th^{-1}[P_1(h)-E_1]\big) (x,y) \Big|^2 dy \, \|u_h(y)\|^2_{L^2(M)} \nonumber\\
=& \int \sum_j \chi(Th^{-1}[\lambda_j^{(1)}(h)-E_1]) B_ku_j^h(x) \overline{u_j^h(y)} \sum_l \chi(Th^{-1}[\lambda_l^{(1)}(h)-E_1]) \overline{B_k u_l^h(x)} u_l^h(y) dy \nonumber\\
=& \sum_j \rho(Th^{-1}[\lambda_j^{(1)}(h)-E_1]) B_k u_j^h(x) \overline{B_k u_j^h(x)} \nonumber\\
=& T^{-1} \int \hat\rho(t/T) (B_k e^{\frac{i}h t(P_1-E_1)} B_k^*)(t,x,x) dt,
\end{align}
with setting $\rho(t)=(\chi(t))^2$.


Now one can insert a new cut-off function $\hat\beta(t)\in C_0^\infty(\mathbb{R})$ satisfying $\hat\beta(t)=1$, $|t|<1$ and $\hat\beta(t)=0$, $|t|>2$ to get that
\begin{align}
\Big| B_k u_h (x) \Big|^2 &\leq  T^{-1} \int \big(1-\hat\beta(t)\big)\hat\rho(t/T) (B_k e^{\frac{i}h t(P_1-E_1)}B_k^*)(t,x,x) dt \nonumber\\
&\qquad\qquad + T^{-1} \int \hat\beta(t)\hat\rho(t/T) (B_k e^{\frac{i}h t(P_1-E_1)}B_k^*)(t,x,x) dt \nonumber\\
&=:I(x,x)+II(x,x).
\end{align}

Following \cite{STZ11} one has that
\begin{equation}\label{BI}
I(x,x) =O(h^\infty).
\end{equation}


For the second part, with setting $\hat\alpha(t)=\hat\beta(t)\hat\rho(t/T)$, one notices that
\begin{align*}
II(x,x)&=T^{-1} \int \hat\beta(t)\hat\rho(t/T) (B_ke^{\frac{i}h t(P_1-E_1)} B_k^*)(t,x,x) dt \\
&=T^{-1} \sum_j \alpha(h^{-1}[\lambda_j^{(1)}(h)-E_1]) B_k u_j^h(x) \overline{B_k u_j^h(x)}.
%&=T^{-1} \sum_j \alpha(h^{-1}[\lambda_j^{(1)}(h)-E_1]) B_k \Big(\alpha \big(h^{-1}[P_2(h)-E_2(h)]\big) u_j^h\Big)(x) \overline{B_k u_j^h(x)}\\
%&=T^{-1} \sum_j \alpha(h^{-1}[\lambda_j^{(1)}(h)-E_1]) B_k \Big(\alpha \big(h^{-1}[\lambda_j^{(2)}(h)-E_2(h)]\big) u_j^h(x) \Big)) \overline{B_k u_j^h(x)} \\
%&=T^{-1} \sum_j \alpha(h^{-1}[\lambda_j^{(1)}(h)-E_1])\alpha \big(h^{-1}[\lambda_j^{(2)}(h)-E_2(h)]\big) B_k u_j^h(x) \overline{B_k u_j^h(x)} \\
%&=  \frac{1}{2\pi} \int \hat\alpha(t) \Big(A e^{\frac{i}h t(P_1-E_1)} \alpha\big(h^{-1}[P_2(h)-E_2(h)]\big)A^* \Big)(t,x,x) dt,
\end{align*}

Following the same argument of \eqref{upshot} and \eqref{sum to integral}, one can apply Lemma \ref{lwl} to get that
\begin{equation}\label{BII}
II(x,x) \leq C T^{-1} I_h(x)^2.
\end{equation}
Combining \eqref{terms}, \eqref{BI} and \eqref{BII} with  noticing that there are $O(\delta^{-(n-1)})$ terms, one can take $T$ sufficient large to get that
\begin{equation}
    |B(x, hD)u_h(x)|=o(1)I_n(h).
\end{equation}



%Note that
%\begin{align}
%&\sum_j \rho_1(h^{-1}[\lambda_j^{(1)}(h)-E_1(h)]) u^h_j(x) \overline{u^h_j(y)} \nonumber \\
%=&\sum_j \rho_1(h^{-1}[\lambda_j^{(1)}(h)-E_1(h)]) \big(\rho_2\big(h^{-1}[P_2(h)-E_2(h)]\big) u^h_j\big) (x) \overline{u^h_j(y)}.
%\end{align}

%Again, by orthogonality of $\{u_j^h\}$, one has that
%\begin{equation}
%\big(\rho_2\big(h^{-1}[P_2(h)-E_2(h)]\big) u^h_j\big) (x)=\rho_2(h^{-1}[\lambda_j^{(2)}(h)-E_2(h)]) u_j^h(x).
%\end{equation}
%Hence we can get 
%\begin{align}
%&\sum_j \rho_1(h^{-1}[\lambda_j^{(1)}(h)-E_1(h)]) u^h_j(x) \overline{u^h_j(y)} \nonumber \\
%=&\sum_j \rho_1(h^{-1}[\lambda_j^{(1)}(h)-E_1(h)])\rho_2(h^{-1}[\lambda_j^{(2)}(h)-E_2(h)])  u^h_j (x) \overline{u^h_j(y)} \nonumber \\
%=&\int_{\mathbb{R}} \hat{\rho_1}(t) \big( \rho_2 \big(h^{-1}[P_2(h)-E_2(h)]\big) e^{\frac{i}h tP_1} \big)(t,x,y) dt.
%\end{align}
%Therefore we need to show that
%\begin{align}
%\sum_j &\rho_1(h^{-1}[\lambda_j^{(1)}(h)-E_1(h)])\rho_2(h^{-1}[\lambda_j^{(2)}(h)-E_2(h)])  \int_{-\frac{1}2}^{\frac{1}2}\int_{-\frac{1}2}^{\frac{1}2}  u^h_j (\tau_1) \overline{u^h_j(\tau_2)}  d\tau_1 d\tau_2 \nonumber \\
%&\leq C_\gamma h.
%\end{align}

%We claim that we need to prove that
%\begin{equation}
%\sum_j \rho_1(h^{-1}[\lambda_j^{(1)}(h)-E_1(h)])\rho_2(h^{-1}[\lambda_j^{(2)}(h)-E_2(h)])  \int_{-\frac{1}2}^{\frac{1}2}\int_{-\frac{1}2}^{\frac{1}2}  u^h_j (\tau_1) \overline{u^h_j(\tau_2)}  d\tau_1 d\tau_2 
%\leq C_\gamma h.
%\end{equation}

%Therefore we need to show that
%\begin{equation}
%\int\int\int_{\mathbb{R}} \hat{\rho_1}(t) \big( \rho_2(P_2/h) e^{\frac{i}h tP_1} \big)(t, \tau_1, \tau_2) dt d\tau_1 d\tau_2
%\end{equation}







%Let us set
%\begin{align}
%II_1&=\int \int   \tilde{I}_1 (\gamma(\tau_1), \gamma(\tau_2); h) \Big( 1-\chi(h^{-1+\delta_0}|\gamma(\tau_1)-\gamma(\tau_2)|) \Big) d\tau_1 d\tau_2, \\
%II_2&=\int \int   \tilde{I}_2 (\gamma(\tau_1), \gamma(\tau_2); h) \Big( 1-\chi(h^{-1+\delta_0}|\gamma(\tau_1)-\gamma(\tau_2)|) \Big) d\tau_1 d\tau_2.
%\end{align}


%Firstly, with taking $\epsilon_0$ sufficiently small, we can find that 
%\begin{equation}
%|\partial_\Theta \Phi(s,x,y,\Theta)|\geq C|x-y|.
%\end{equation}
%Hence one can show that
%\begin{equation}
%| II_1 |=O(h^\infty)
%\end{equation}
%with using integration by parts in variable $\Theta$ since each integration will gain $O(h^{\delta_0})$.

%Finally, we need to deal with the term $II_2$. One applies the stationary phase theorem in variable $\Theta$ with \cs\cs the non-degenerate condition of $p_2$. One has
%\begin{align}
%\tilde{I}_2 (\gamma(\tau_1), \gamma(\tau_2); h),
%\end{align}

% the possible $\Theta$-critical points of the phase $\Phi(s,x,y,\Theta)$ must satisfy
%\begin{align}
%x-y+s\partial_\Theta p_2(x,\Theta)+O(s^2) =0 \label{critical2}
%\end{align}

%If $s\leq Ch$ for some large constant $C$. Let us take a new cut-off function $\chi\in C_0^\infty(\mathbb{R})$ with $\chi(\tau)=1$ if $\tau\in[0, 1]$ and $0\leq\chi(\tau)\leq1$. One can deduce that
%\begin{align}\label{small s1}
%\Big| h^{-1}\int\int \int_{0\leq s\leq Ch}\int \hat{\rho_2}(s)  e^{\frac{i}h \Phi(s, \gamma(\tau_1), \gamma(\tau_2), \Theta)} \chi(h^{-\frac{1}2}|\gamma(\tau_1)-\gamma(\tau_2)|)
% c'()d\Theta ds d\tau_1d\tau_2 \Big|\leq C_M h^{1/2}.
%\end{align}
%On the other hand one can show that
%\begin{align}\label{small s2}
%|h^{-1}\int\int \int_{0\leq s\leq Ch}\int \hat{\rho_2}(s)  e^{\frac{i}h \Phi(s, \gamma(\tau_1), \gamma(\tau_2), \Theta)} \big(1-\chi(h^{-\frac{1}2}|\gamma(\tau_1)-\gamma(\tau_2)|) \big) \nonumber\\
% c'()d\Theta ds d\tau_1d\tau_2 \Big|=O(h^\infty)
%\end{align}
%with using integration by parts in variable $\Theta$ since each integration will gain $O(h^{1/2})$.
 %one has the nontrivial contribution in \eqref{upshot1} is \cs\cs $|x-y|\leq C'h$ and $C'h\leq |x-y|\leq C'h^{1/2}$where $|\gamma(\tau_1)-\gamma(\tau_2)|$ is comparable to $h$ due to \eqref{critical2} and the generic assumption of $\gamma$. Hence one can deduce that
%\begin{equation}\label{small s}
%\Big| h^{-1}\int\int \int_{0\leq s\leq Ch}\int \hat{\rho_2}(s)  e^{\frac{i}h \Phi(s, \gamma(\tau_1), \gamma(\tau_2), \Theta)}
% c'()d\Theta ds d\tau_1d\tau_2 \Big|\leq C_M h 
%\end{equation}
%which is better than \eqref{upshot}.

%\cs\cs\cs\cs\cs\cs\cs\cs\cs

%If $|x-y|<h^{1-\delta_0}$, here $\delta_0$ is a sufficiently small positive constant, one can apply stationary phase in $(\Theta,y)$ to get $h^{1-\delta_0}$.

%If $|x-y|\geq h^{1-\delta_0}$.
%When $s<|x-y|$, one can apply non-stationary phase.
%When $s\geq |x-y|$, with the Morse condition of $p_2$, one can apply stationary phase in $\Theta$. Then integrate in $s$. Finally one can use non-stationary phase in $|x-y|$.

%\cs\cs\cs\cs\cs\cs\cs\cs\cs



%\cs\cs\cs\cs\cs\cs\cs\cs\cs

%If $s\geq Ch$, one has that the Hessian $\frac{\partial^2\Phi}{\partial(s,\Theta)}$ is \cs\cs not quite right
%\begin{equation}\label{Hessian}
%\begin{pmatrix}
%O(1) & \partial_\Theta p_2+ O(s) \\
%\partial_\Theta p_2+ O(s)  & \partial^2_\Theta p_2\cdot O(s)
%\end{pmatrix}.
%\end{equation}

%If $|\partial_\Theta p_2|\geq Cs^{\frac{1}2}$ for some large constant $C$, the Hessian \eqref{Hessian} is non-degenerate. One can apply stationary phase in $(\Theta, s)$ to get
%\begin{equation}
%\tilde{I} (\gamma(\tau_1),\gamma(\tau_2); h)= \exp[\frac{i}h \Phi(\gamma(\tau_1), \gamma(\tau_2))]  \cdot c'(),
%\end{equation}
%here $\Phi(\gamma(\tau_1), \gamma(\tau_2))=\left<\gamma(\tau_1)-\gamma(\tau_2), \Theta_0\right>+O(s_0^2)$ where $s_0(\tau_1,\tau_2)$ and $\Theta_0(\tau_1,\tau_2)$ is the non-degenerate critical point solving \eqref{critical1} and \eqref{critical2}, and 
%\begin{equation}
%|c'()|\leq C's_0^{-1/2}.
%\end{equation}

%If $|\partial_\Theta p_2|\leq Cs^{\frac{1}2}$, the Hessian \eqref{Hessian} may be non-degenerate.
%If $|x-y|\geq C'h^{\frac{1}2}$ with another large constant $C'$, one can use integrate by parts.
%If $|x-y|<C'h^{\frac{1}2}$, one can use stationary phase in $(\Theta, y)$ to get



%\cs\cs\cs\cs\cs\cs\cs\cs\cs

%If $s\geq Ch$,  one has that the Hessian $\frac{\partial^2\Phi}{\partial(s,\Theta)}$ is
%$$\begin{pmatrix}
%O(1) & \partial_\Theta p_2+ O(s) \\
%\partial_\Theta p_2+ O(s)  & \partial^2_\Theta p_2\cdot O(s)
%\end{pmatrix}.$$
%Due to the admissible assumption of $p_2$ with taking the $\epsilon$ small enough in \eqref{spt rho2}, one can find that the Hessian $\frac{\partial^2\Phi}{\partial(s,\Theta)}$ is non-degenerate. Hence the $(\Theta, s)$-critical points of the phase $\Phi(s,x,y,\Theta)$ are isolated and so, without loss of generality, we assume that there is only one.
%With choosing a large constant $C$ and applying stationary phase in $(\Theta, s)$, one has
%\begin{equation}
%\tilde{I} (\gamma(\tau_1),\gamma(\tau_2); h)= \exp[\frac{i}h \Phi(\gamma(\tau_1), \gamma(\tau_2))]  \cdot c'(),
%\end{equation}
%here $\Phi(\gamma(\tau_1), \gamma(\tau_2))=\left<\gamma(\tau_1)-\gamma(\tau_2), \Theta_0\right>+O(s_0^2)$ where $s_0(\tau_1,\tau_2)$ and $\Theta_0(\tau_1,\tau_2)$ is the non-degenerate critical point solving \eqref{critical1} and \eqref{critical2}, and 
%\begin{equation}
%|c'()|\leq C's_0^{-1/2}.
%\end{equation}

%\cs\cs Maybe this is better to give a geometric interpretation of $\Phi(\gamma(\tau_1), \gamma(\tau_2))$ which is like a distance function under the flow induced by $p_2$.\cs\cs
%Finally, we can integral in variables $(\tau_1, \tau_2)$ with reminding that the nontrivial part of \eqref{upshot1} is where $|\gamma(\tau_1)-\gamma(\tau_2)|$ is comparable to $h$ due to \eqref{critical2} to get that
%\begin{equation}\label{large s}
%\Big|\int\int\tilde{I} (\gamma(\tau_1),\gamma(\tau_2); h) d\tau_1 d\tau_2\Big|= C_M h^{1/2}.
%\end{equation}
%Combining \eqref{small s1}, \eqref{small s2} and \eqref{large s} one has \eqref{upshot}.






%We also know the wave operator $e^{it (-h^2\Delta_g - 1) }$ has a Schwartz kernel for small time $t$,
%\begin{equation}
%K(t,x,y,h)=(2\pi h)^{-n}\int_{\mathbb{R}^n} e^{i[\phi(t,x,\xi)-\big<y,\xi\big>]/h} a(t,x,y,\xi;h) d\xi +O_{C^\infty}(h^\infty),
%\end{equation}
%with\cs\cs $a\sim $ and $\phi(t,x,\xi)$ solves the eikonal equation
%\begin{equation}
%\end{equation}
%So after taking $\delta$ small enough so that the above paramatrix is valid, we reduce to proving that
%\begin{equation}
%\Big| \int_{t\leq \delta} \hat \chi(t/T) K(t,x,x,h) dt \Big| = O(h^{1-n}).
%\end{equation}








%\cs\cs On negatively curved surface, following SZ method carefully you can see $o(1)$ improvement. Notice why you can't ge $o(1)$ improvement when the curvature is zero, the main fact you need the volume factor from Gunther's comparison theorem to control extra term. Indeed when taking summation, you can't get $T^{-\epsilon}$ improvement.

%Let $\Pi$ denote the space of all unit-length geodesics in $M$, \cs\cs say something why you focus on $\Pi$ or something periodic curve?, then our main result is the following $o(1)$ improved `non-concentration' theorem,  \cs\cs $L^p$ estimate and log improvement (c.f. Blair Soge 2018 JDG)

%\begin{theo}\label{main}
%Assume that $(M,g)$ is as above. There is an $h_0>0$ so that
%\begin{equation}\label{main eq}
%\sup_{\gamma\in \Pi}\| h \partial_\nu u_{h} \|_{L^2(\gamma)}=O(-\log h)^{-1/4}, \quad h<h_0.
%\end{equation}
%\end{theo} \cs\cs the power $-\frac{1}4$ may be wrong. I guess it may be $-\frac{1}2$. Get an $L^p$ estimate.







%For the terms coming from $\partial_{\nu_x}$ and $\partial_{\nu_y}$ acting on $a(x,y,\xi)$, following \cite{BGT}(or following proof), one can easily show the corresponding operator's $L^2$ norm is bounded by $O(\lambda^{1/2})$ which is favourable. So we only take care of the terms coming from $\partial_{\nu_x} \partial_{\nu_y} e^{i\psi(x,y)\cdot \xi \pm it|\xi|}$.





\bigskip


\section{Examples}\label{section4}

\subsection{Surface of revolutions}
We study the surface of revolution as an example to demonstrate several points mentioned in the Introduction. Proposition \ref{SOR degen} shows the geometric meaning of the rank $1$ condition in this simple case; Proposition \ref{SOR recur} and \ref{SOR quasi} clarifies the point made in Remark (1) of Theorem \ref{main1}. We believe that some constructions can be generalized to Riemannian warped products.

In general, a surface of revolution embedded in $\R^3$ is given by parametrization
\begin{equation}\label{SoR parametrization}
    \begin{cases}
        x=f(t)\cos\phi\\
        y=f(t)\sin\phi\\
        z=g(t),
    \end{cases}
\end{equation}
for $t\in [t_-,t_+]$ and $\phi\in [0,2\pi]$. Its first fundamental form is 
\begin{align*}
    E&=e_t\cdot e_t=f'(t)^2+g'(t)^2\\
    F&=e_t\cdot e_\phi=0\\
    G&=e_\phi\cdot e_\phi=f^2(t)
\end{align*}
and we want to arrange for $f'(t)^2+g'(t)^2=1$, so that the metric is $ds^2=dt^2+f^2(t)d\phi^2$. We call this metric $g$, not to be confused with the above function $g$. One can compute the Laplacian in the coordinate system $(t,\phi)$,
\begin{align}
    \Delta &=\frac{1}{\sqrt{|g|}}\sum\partial_j(\sqrt{|g|}\,g^{jk}\partial_k)\nonumber\\
    &=\partial_t^2+f^{-2}(t)\partial_\phi^2+\frac{f'(t)}{f(t)}\partial_t.
\end{align}
The principal symbol is then $|\xi|^2_g=\xi_t^2+f^{-2}(t)\xi_{\phi}^2$. The quantum integrable system is defined by $\mathcal{P}=(p_1,p_2)$, where $p_1=|\xi|^2_g$, $p_2=\xi_{\phi}$.

We will make extra assumptions in order to make the problem simpler. We choose $f$ such that $f(t_-)=f(t_+)=0$ and assume that $f$ makes the metric smooth at the poles. We also assume that $f''(t_0)\neq 0$ for the $t_0$ where $f'(t_0)=0$. The assumption guarantees that the set $\{(t,\phi)\in M:f'(t)=0\}$ consists of disjoint union of horizontal circles. For simplicity, we may assume there is only one such circle, which we refer to as the \textbf{equator}.

There is a simple characterization of the rank $1$ condition in this case.

\begin{prop}\label{SOR degen}
    Let $M$ be a surface of revolution as above, and $x\in M$ not a pole nor on the equator. Then $\mathcal{P}$ satisfies the rank $1$ condition on $\Gamma=\{(\xi_t,\xi_{\phi})\in S^*_xM: \xi_t\neq 0\}$ at $x$.
\end{prop}

\begin{proof}
    At $x=(t,\phi)$, the coshpere is an ellipse $\mathcal{C}_x=\{\xi:\xi_t^2+\frac{\xi_{\phi}^2}{f^2(t)}=1\}$. Then, ${p_2}|_{\mathcal{C}_x}=\xi_{\phi}|_{\mathcal{C}_x}$ is critical only if $\xi_{t}=0$, that is, when the covector is vertical. Then, $\mathcal{P}$ satisfies the rank $1$ condition on $\{\xi_t\neq 0\}$ at $x$.  
\end{proof}

\begin{rem}
    Note that the trick in Section \ref{section superint} cannot be used here.
\end{rem}

The next proposition characterizes the recurrent set of certain surfaces of revolution.

\begin{prop}\label{SOR recur}
    Let $M$ be a surface of revolution as above with a real-analytic profile function $f$, and that $g$ is not a Zoll metric. Let $x\in M$ be a point on the equator, then the recurrent set $\mathcal{R}_x$ is a dense subset in $S^*_xM$ and $\mu(\mathcal{R}_x)=0$, where $\mu$ is the induced Liouville measure on $S^*_xM$.
\end{prop}

\begin{rem}
    In this case, $\mathcal{R}_x=\widetilde{\mathcal{L}}_x=\mathcal{L}_x$. Also, the conclusion holds actually for any $x$ that is not a pole. To keep the proof short, we assume that it lies on the equator. 
    
    We can almost produce a one-line proof following \cite[Thm 3.5.24]{Kli95}, which claims that the geodesic flow on a surface of revolution is "equivalent" to certain linear torus action, the slope of which is explicit. However, we opt for an elementary proof as below.
\end{rem}

\begin{figure}[h]
\includegraphics[width=\textwidth]{drawing2.pdf}
\caption{The relation of the equator (red) and the geodesic with initial angle $\psi$ (blue) in the proof. 
Left: trajectories on the surface. Middle: canonical coordinates $(t,\phi)$. Right:  geodesic normal coordinates $(s,\psi)$.}
\label{Figure 2}
\end{figure}


\begin{proof}
We invoke the Clairaut relation that holds on every surface of revolution: For a geodesic $l$, the quantity $\gamma=f(t)\sin(\psi(t))$ is a constant along $l$, where $\psi(t)$ is the angle between the geodesic $l$ and the meridian it intersects with at $(t,\phi)$. Write $\psi=\psi(t_0)$ for the initial angle. It is known that the geodesic oscillates within a tubular region near the equator (\cite{Kli95}).

Fix $x$ on the equator. Without loss of generality, let $\phi(x)=0$. For a geodesic $l$ starting at $x$ with angle $\psi$ to the meridian, it is parametrized by arc length so that $l(s)=(t(s),\phi(s))$, $t'(s)^2+f'(t(s))^2\phi'(s)^2=1$. It satisfies the geodesic equations
\begin{equation}
    \begin{cases}
        \frac{d^2t}{ds^2}=f(t(s))\,\frac{df}{ds}\left(\frac{d\phi}{ds}\right)^2\\
        \frac{d}{ds}\left(f^2(t(s))\frac{d\phi}{ds}\right)=0
    \end{cases}
\end{equation}
subject to the initial conditions
\begin{align}
    \phi(0)=0,\quad t(0)=0,\quad \frac{d\phi}{ds}(0)=\frac{\sin\psi}{f(t_0)},\quad \frac{dt}{ds}(0)=\cos\psi.
\end{align}

The only free variable in the initial conditions is $\psi$. Using the fundamental theorem of ODE, one can write a real-analytic function $L:(\psi,s)\mapsto(t(\psi,s),\phi(\psi,s))=(L_1,L_2)$, which serve as a $C^{\omega}$ change-of-coordinates into the geodesic normal coordinates $(\psi,s)$. Then, the equator as a curve $\{(t,\phi):t=0\}$ can be expressed as $\{(\psi,s):s=S(\psi)\}$, a graph of some real-analytic function $S$. Therefore, the function $\Phi(\psi)=L_2(\psi,S(\psi))$ is real-analytic, possibly multi-valued function, which denotes the change of longitude when the geodesic $l$ returns to the equator. We take the smallest positive value as a well-defined branch, and still denote by $\Phi$. Its inverse $\Phi^{-1}$ is also real-analytic in a suitable domain.

The geodesic is recurrent if and only if $\Phi(\psi)$ is a rational multiple of $2\pi$. Let $B:=\{\text{rational multiples of }2\pi\}$. Suppose $\exists A\subset S^1,\, \mu(A)>0$ such that $\Phi(A)= B$. Then $\exists A'\subset A$, $\mu(A')>0$, $\Phi(A')=\{c\}$. By analyticity, $\Phi=const.$ Since $M$ possess rotational symmetry, we get the same constant for every geodesic. It forces the manifold to be a Zoll surface.

By assumtion, the surface of revolution is not Zoll. Then, the inverse image $\Phi^{-1}(B)$ must have zero measure. Denseness follows from the fact that $\Phi^{-1}$ is $C^{\omega}$ and that $B$ is dense.

% $$
% \omega(\gamma)=\int_I\frac{\gamma}{f(t)\sqrt{f^2(t)-\gamma^2}}dt
% $$
% is a rational multiple of $\pi$ if and only if $l$ is closed geodesic. (For the round sphere, $\omega$ is always $0$.)
\end{proof}


We mention a method for constructing a sequence of joint quasimodes from a sequence of eigenfunctions of its horizontal factor. We believe that such a construction also exists in certain Riemannian warped products. 


\begin{prop}\label{SOR quasi}
    Let $M$ be a surface of revolution as above with profile function $f$ such that $f\leq 1$, and $\frac{|f'|}{|1-f^2|}$ is bounded where $f=1$. Then there exists an $O(h)$-quasimode $u_h$ such that $|u_h(t_0,\phi)|\simeq h^{-1/4}$ and $|u_h(t,\phi)|\simeq  e^{-C(t,\phi)/h}$ for $f(t)\neq 1$. 
\end{prop}
\begin{proof}
    Separation of variables: Let $u_{\lambda}(t,\phi)=T(t)e^{i\lambda\phi}$. We call the following \textbf{highest-weight quasimode}: 
\begin{equation}\label{trial}
    u_{\lambda}(t,\phi)=\lambda^{1/4}\exp\{\lambda\int\sqrt{\frac{1-f^2}{f^2}}dt\}\,e^{i\lambda \phi}.
\end{equation}
One can check by direct computation that 
\begin{align}
\Delta u_{\lambda}&=\Big(\partial_t^2+f^{-2}(t)\partial_\phi^2+\frac{f'(t)}{f(t)}\partial_t\Big)u_{\lambda} \nonumber\\
&=\Big(-\lambda^2-\lambda\frac{f'(t)}{f(t)^2\sqrt{1-f(t)^2}}+\lambda\frac{f'(t)\sqrt{1-f(t)^2}}{f(t)^2}\Big)u_{\lambda}\nonumber\\
&=\Big(-\lambda^2-\frac{\lambda f'(t)}{\sqrt{1-f(t)^2}}\Big)u_{\lambda}
\end{align}
At points where $f(t_0)=1$, we expect that on the equator $f'(t_0)=0$ at correct order s.t. the second term does not blow up. It is guaranteed by our assumption.

\end{proof}

\begin{rem}
Back to the round sphere, we let $f=\cos t$, then 
$$
\exp\{\lambda\int\sqrt{\frac{|1-f^2|}{f^2}}dt\}=\exp\{\lambda\int |\tan t|dt\}=\exp\{\lambda \ln|\cos t|\}=(\cos t)^{\lambda}. 
$$
(in this convention, $t\in [-\pi/2,\pi/2]$, so $\cos t>0$.)
\end{rem}

 
Combining the Proposition \ref{SOR recur} and \ref{SOR quasi}, we exemplified Remark (1) of Theorem \ref{main1}, and partly justified the necessity of the rank $k$ condition of nondegeneracy.

Finally, we mention that most results in Section \ref{section4} can be generalized to Riemannian warped products $I \times_fN$. It seems interesting to explore their utilities as (counter-)examples in higher dimensions.

The following examples are inspired by \cite{GT2020}.

\subsection{Laplacians on Liouville tori}

Consider the two-torus $M=\mathbb{R}^2/\mathbb{Z}^2$ with a Liouville metirc given by $g=(a(x_1)+b(x_2))(dx_1^2+dx_2^2)$, here $a, b:\mathbb{R}/\mathbb{Z}\to \mathbb{R}^+$ are two, smooth, positive periodic functions with $\min_{0\leq x_1\leq1} a(x_1)>\max_{0\leq x_2 \leq 1}b(x_2)$. The associated Laplacian $$P_1(h)=-[a(x_1)+b(x_2)]^{-1}((h\partial_{x_1})^2+(h\partial_{x_2})^2)$$ is QCI with $$P_2(h)=-[a(x_1)+b(x_2)]^{-1}(b(x_2)(h\partial_{x_1})^2-a(x_1)(h\partial_{x_2})^2).$$

Taking any point $z_0=(x_0, y_0)\in\mathbb{R}^2/\mathbb{Z}^2$ with $\alpha=a(x_0)>b(y_0)=\beta$ we have that $$p_2\vline_{\,T_{z_0}^*}=\beta(\alpha+\beta)^{-1}\xi^2- \alpha(\alpha+\beta)^{-1}\eta^2.$$
Since $S_{z_0}^*=\{(\xi,\eta);\, \xi^2+\eta^2=\alpha+\beta>0\}$, one sets $\xi=\sqrt{\alpha+\beta}\cos\theta$ and $\eta=\sqrt{\alpha+\beta}\sin\theta$. Hence $$p_2\vline_{\,S_{z_0}^*}(\theta)=\beta\cos^2\theta-\alpha\sin^2\theta,$$
which is critical at $\theta=\theta_k=k\pi/2$, $k=0,1,2,3$. One can find the corresponding energy $E_2^k:=p_2(\theta_k)$. Hence {\em rank 1 condition} is valid if $(1, E_2)\in \mathcal{B}\backslash\{(1,E_2^k)\}$, here $\mathcal{B}$ is the set of regular values of the moment map $\mathcal{P}$.

\subsection{Laplacians on ellipsoids}
Consider the ellipsoid $\mathcal{E}=\{x\in \mathbb{R}^3; \sum_{j=1}^3\frac{x^2_j}{a^2_j}=1\}$ where ${0<a_3<a_2<a_1}$ are fixed constants. It is QCI with reduced Hamiltonian $H$ and the reduced integral in involution $P$ which are given by the formulas \cite{Tot97}:
\begin{align*}
H&=a_3(x_1\xi_2-\xi_1x_2)^2+a_2(x_1\xi_3-x_3\xi_1)^2+a_1(x_3\xi_2-x_2\xi_3)^2,\\
P&=(x_1\xi_2-\xi_1x_2)^2+(x_1\xi_3-x_3\xi_1)^2+(x_3\xi_2-x_2\xi_3)^2.
\end{align*}
Here $P$ is equal to the principle symbol of the standard Laplacian on $S^2$.
Firstly, we consider the point  $(x_0, \xi_0)$  on the hyperbolic geodesics $\Gamma^{\pm}=\{(x,\xi)\in S^*(S^2); x_2=\xi_2=0\}$. One has that
\begin{align*}
H\vline_{\,\Gamma^\pm}&=a_2(x_1\xi_3-x_3\xi_1)^2,\\
P\vline_{\,\Gamma^\pm}&=(x_1\xi_3-x_3\xi_1)^2=\frac{1}{a_2}  H |_{\,\Gamma^\pm}.
\end{align*}
Hence {\em rank k condition} fails. As \cite{Tot97} points out that one can find a sequence of eigenfunctions $\{u_h\}$ such that $\|u_h\|_{L^\infty(\Gamma(h))}=O(h^{1/2}/|\ln h|)$, here $\Gamma(h)$ is a tube of width $O(h)$ about a hyperbolic geodesic $\Gamma$.


Now we consider the point $x_0$ away $\pi(\Gamma)$ where $\pi$ is the standard projection map. Denote the rectangles $R_+=(0, T_1)\times (0, T_2)$ and $R_-=(T_1, 2T_2)\times(0, T_2)$. We let $\Phi_\pm: R_\pm\to \mathcal{E}\cap \{w_2\neq 0\}$ be the conformal mapping sending vertices of $R_\pm$ to the four umbilic points $p_k; k=1,\dots 4$ of $\mathcal{E}$. We choose orientations so that $\Phi_+(x,T_2)=\Phi_-(2T_1-x,T_2)$ and $\Phi_+(x,0)=\Phi_-(2T_1-x,0)$. We henceforth let $\Phi:=\Phi_\pm: R\to \mathcal{E}$ denote the induced conformal mapping with $\Phi\vline_{R_\pm}=\Phi_{\pm}$ and $R:=R_+\cup R_-$. One has that the intrinsic Riemannian metric $\mathcal{E}$ pulled-back to $R$ is locally of Liouville form \cite{GT2020} $$ds^2=\big(a(x_1)+b(x_2)\big)(dx_1^2+dx_2^2).$$

Then the same argument works as in the Liouville torus. For the specific joint eigenvalues $\{(1, E_2^k)\}$ whose corresponding geodesics including ones that connect $x_0$ and the four umbilic points and have a focusing effect which are forward and backwards asymptotic to the middle-length ellipse \cite{Tot97, GT2020}, the {\em rank k condition fails}. However {\em rank 1 condition} is valid if $(1, E_2)\in \mathcal{B}\backslash\{(1,E_2^k)\}$ which is consistant with the result of \cite{Tot97} which shows that there exist a sequence of eigenfunctions $\{u_h\}$ such that $\|u_h\|_{L^\infty}=O(1/|\ln h|)$ outside an arbitrarily small (but fixed) neighbourhood of $\Gamma$.

% \subsubsection{Biaxial ellipsoid} We consider the concrete example 
% $$
% M=\{(x,y,z)\in \R^3:\frac{x^2}{a^2}+\frac{y^2}{a^2}+\frac{z^2}{b^2}=1\}.
% $$
% Let $\beta(t)$ be the inverse function of the elliptic integral
% \begin{equation}\label{elliptic int}
%     t(\beta)=\int_0^{\beta}\sqrt{a^2\sin^2 s+b^2\cos^2s}\,ds.
% \end{equation}
% We claim that $f(t)=a\cos(\beta(t))$, $g(t)=b\sin(\beta(t))$ satisfy $f'(t)^2+g'(t)^2=1$. In fact, $t'(\beta)=\sqrt{a^2\sin^2 \beta+b^2\cos^2\beta}$. By inverse differentiation,
% $$
% \beta'(t_0)=\frac{1}{t'(\beta(t_0))}=\frac{1}{\sqrt{a^2\sin^2 \beta(t_0)+b^2\cos^2\beta(t_0)}}.
% $$
% Therefore,
% \begin{align}
%     &f'(t)^2+g'(t)^2\nonumber\\
%     =&a^2\beta'(t)^2\sin^2(\beta(t))+b^2\beta'(t)^2\cos^2(\beta(t))\nonumber\\
%     =&\Big(a^2\sin^2(\beta(t))+b^2\cos^2(\beta(t))\Big)\beta'(t)^2=1.
% \end{align}
% The metric is
% \begin{equation}\label{FFF}
%     ds^2=dt^2+a^2\cos^2\beta(t)d\phi^2.
% \end{equation}
% Furthermore, we can compute the Laplacian in $(t,\phi)$ coordinates,
% \begin{align}
%     \Delta &=\frac{1}{\sqrt{|g|}}\sum\partial_j(\sqrt{|g|}\,g^{jk}\partial_k)\nonumber\\
%     &=\partial_t^2+f^{-2}(t)\partial_\phi^2+\frac{f'(t)}{f(t)}\partial_t\nonumber\\
%     &=\partial_t^2+\frac{1}{a^2\cos^2\beta(t)}\partial_\phi^2-\frac{\tan(\beta(t))}{\sqrt{a^2\sin^2\beta(t)+b^2\cos^2\beta(t)}}\partial_t.
% \end{align}
% The last term is of lower order ($O(h)$ negligible).

% \cs\cs\cs

% One parametrizes convex surfaces of revolution by using geodesic polar coordinates $(t, \varphi)\in [-1,1]\times [0, 2\pi]$ in terms of which
% \begin{equation*}
% p_1(t,\varphi; \xi_t, \xi_\varphi)=\xi_t^2+f^{-2}(t)\xi_\varphi^2,
% \end{equation*}
% and
% \begin{equation*}
% p_2(t,\varphi; \xi_t, \xi_\varphi)=\xi_\varphi.
% \end{equation*}
% The profile function satisfies $f(-1)=f(1)=0$ and $f(t)$ is a non-negative Morse function with a single non-degenerate maximum at $t=t_0\in(-1,1)$.

% \cs\cs\cs

% {\bf Case 1.}  If one can use $\xi_\varphi$ to parametrize $p_2$. Hence
% \begin{equation}
% \partial_{\xi_\varphi} p_2=1,
% \end{equation}
% which is never zero.

% {\bf Case 2.} If one can use $\xi_t$ to parametrize $p_2$. Hence
% \begin{equation}
% \partial_{\xi_t} p_2^2=f^2(t) \partial_{\xi_1}(1-\xi_t^2)=-2f^2(t) \xi_t.
% \end{equation}
% If $\xi_t\neq 0$,

% If $\xi_t=0$. When $t\neq t_0$, the point $x_0$ is allowable since $\xi_\varphi\neq 0$ in $\mathcal{R}_x$, here $x=(t,\varphi)$ \cs\cs. but when $t= t_0$, the point $x_0$ is not allowable.
% For an example, set $f(t)=\sqrt{1-t^2}$ near $t_0=0$ which means that near $t_0=0$ the surface is just a sphere.

% \cs\cs\cs

% Consider the point $x=(t_1, \varphi)$, where $t_1\neq \pm1$. One has that
% \begin{equation*}
% \mathcal{C}_x=\{(t_1,0; \xi_t, \xi_\varphi);\, \xi_t^2+f^{-2}(t_1)\xi_\varphi^2=1\}.
% \end{equation*}
% On the set $\mathcal{C}_x\cap\{ |\xi_\varphi|\leq \delta\}$ where $\delta>0$ is sufficiently small, one has that $\partial_{\xi_t} p_1=2\xi_t$ which is away from zero. Hence, one can use $\xi_\varphi$ to parametrize $\mathcal{C}_x$:
% $$
% \mathcal{C}_x\cap\{ |\xi_\varphi|\leq \delta\} =\{(t_1,0,F(\xi_{\phi}),\xi_{\phi}): |\xi_\varphi|\leq \delta\}.
% $$
% Now by the chain rule,
% \begin{equation}\label{near 0}
% \partial (p_2|_{\mathcal{C}_x\cap\{ |\xi_\varphi|\leq \delta\}})=F'(\xi_{\phi})\partial_{\xi_t} p_2+\partial_{\xi_{\phi}}p_2=1\neq0.
% \end{equation}
% which means that \eqref{C1} is valid. 

% On the other hand, on the set $\mathcal{C}_x\cap\{ |\xi_\varphi-E_2|< \delta\}$ where $E_2\neq0$ and $\delta>0$ is sufficiently small, one has that $\partial_{\xi_\varphi} p_1=2f^{-2}(t_1) \xi_\varphi\neq 0$. Hence one can use $\xi_t$ to parametrize $\mathcal{C}_x$:
% $$
% \mathcal{C}_x\cap\{|\xi_{\phi}-E_2|<\delta\}=\{(t_1,0,\xi_t,G(\xi_t)):|G(\xi_t)-E_2|<\delta\},\quad G(\xi_t)=\sqrt{(1-\xi_t^2)f(t_1)^2}
% $$
% Now by the chain rule, 
% \begin{equation}\label{away 0}
% \partial (p_2|_{\mathcal{C}_x\cap\{|\xi_{\phi}-E|<\delta\}})=\partial_{\xi_t}p_2+G'(\xi_{t})\partial_{\xi_{\phi}}p_2
% =G'(\xi_t)=\frac{-f(t_1)^2\xi_t}{\sqrt{(1-\xi_t^2)f(t_1)^2}}.
% \end{equation}
% Then, apart from the poles (on which $f(t_1)=0$), (\ref{away 0}) vanish only when $\xi_t=0$, that is, when $\xi_{\phi}^2=f(t_1)^2$. Hence \eqref{C1} is valid if $f(t_1)\neq \pm E_2$. 

% In fact, this condition also incorporates the previous case when $E_2=0$, because the condition $f(t_1)\neq 0$ excludes exactly the two poles.

% \begin{prop}
%     For a surface of revolution defined as above, $\mathcal{P}$ is non-critical for all $x$ but those satisfying $f(t(x))\neq \pm E_2$. 
% \end{prop}
% \begin{cor}
%     For a surface of revolution defined as above, a sequence of joint eigenfunctions $u_h$ satisfy the bound $|u_h(x)|\leq C h^{\frac{2-n}{2}}|\ln h|^{\frac{1}{2}}$ for those $x$ with $f(t(x))\neq \pm E_2$.
% \end{cor}


% Next we focus on the case of standard sphere $\mathbb{S}^2$. In this case, $f(t)=\sqrt{1-t^2}$ and $t_0=0$.

% Now we use the longitudinal coordinate $\theta\in [0, \pi]$ and latitudinal coordinates $\phi\in [0, 2\pi)$ so that $\mathbb{S}^2\owns x=(\sin\theta \cos \phi, \sin \theta \sin \varphi, \cos \theta)$. One takes the zonal harmonics
% $$u_k=C_k P_k(\cos\theta)$$
% which solves
% $$-\Delta u_k= k(k+1) u_k,$$
% here $C_k\sim k^{\frac{1}2}$ is the $L^2$ normalisation factor, and $P_k$ is the associated Legendre polynomial. The eigenfrequency $h^{-1}=\sqrt{k(k+1)}\sim k$. In this case, the corresponding $h$-Laplacian $P_1(h):=-h^2\Delta$ with eigenvalue $E_1(h)=1+O(h)$ is QCI with commuting $P_2(h)=hD_\varphi$ with $E_2(h)=0$.

% It is well known that $u_k$ saturates the $O(h^{-1/2})$ in an $O(h)$-neighborhood of the poles.

% Next consider a point $x=(\sin\theta_0 \cos \phi, \sin \theta_0 \sin \varphi, \cos \theta_0)$, here $\theta_0\neq0$ and $\theta_0\neq \pi$, which is outside the $O(h)$-neighborhood of the poles. From above argument \eqref{near 0}, we know that \eqref{C1} is valid. From \cite[Theorem 8.21.2]{Sze}, one has
% $$P_k(\cos\theta_0)=\sqrt{\frac{2}{\pi k \sin \theta_0}}\cos\big[(k+\frac{1}2)\theta_0-\frac{\pi}4\big]+O(k^{-3/2}).$$
% One can easily get that
% $$|u_k| \leq C.$$
% This bound is consistent with (and slightly stronger than) the general $O(|\ln h|^{1/2})$ bound given in Theorem \ref{main1}.

% \cs There should be an intermediate case saturating a pair of latitude circles of $0<E_2<1$.\cs

% Next we consider the highest-weight spherical harmonics
% \begin{equation*}
% u_\lambda(x)=\lambda^{1/4}(x_1+i x_2)^{\lambda}=\lambda^{1/4}[\sin \theta \cos \varphi+i \sin \theta \sin \varphi]^\lambda=\lambda^{1/4} (\sin\theta)^\lambda e^{i\lambda \varphi},
% \end{equation*}
% where $\int_M | u_\lambda|^2 d\text{Vol}\sim 1$, and the eigenfrequency $h^{-1}=\lambda=n;\,n=1,\,2,\,3,\dots$. 
% In this case, the corresponding $h$-Laplacian $P_1(h):=-h^2\Delta$ is QCI with commuting $P_2(h)=hD_\varphi$ with $E_2(h)=1$.

% At point $x=(\sin\theta_0 \cos \phi, \sin \theta_0 \sin \varphi, \cos \theta_0)$, the eigenfunction
% \begin{equation*}
% u_\lambda(x)=\lambda^{1/4} (\sin\theta_0)^\lambda e^{i\lambda \varphi},
% \end{equation*}
% here $\theta_0\neq0$ and $\theta_0\neq \pi$. In order to make \eqref{C1} be valid, $\theta_0\neq \frac{\pi}2$ due to above argument \eqref{away 0}. Hence
% \begin{equation*}
% | u_\lambda(x)  | = \big| \lambda^{1/4}  (\sin \theta_0)^{\lambda} \big|  \ll |\ln \lambda|^{\frac{1}2}.
% \end{equation*}
% This bound is also consistent with (and much stronger than) the general $O(|\ln h|^{1/2})$ bound given in Theorem \ref{main1}.


% \subsection{Riemannian warped product}
% Consider Riemannian warped products. It should in some sense resemble the case of surface of revolution.

% Let $g=dt^2+f(t)^2h$, and $(M,g)=(I,dt)\times_f(N,h)$. (I guess $|\xi|_g^2=\xi_t^2+f^{-2}(t)|\xi_\phi|^2_{h}$. For a point $x\in M$ we can write $x=(t(x),\phi(x))$ or just $(t,\phi)$ for short.
% \begin{align}
% p_1(t,\phi,\xi_t,\xi_{\phi})&=\xi_t^2+f^{-2}(t)|\xi_\phi|^2_{h}\\
% p_2(t,\phi,\xi_t,\xi_{\phi})&=|\xi_{\phi}|^2_h
% \end{align}
% It is straight forward to check that $\{p_1,p_2\}=0$.
% A minor question is to make $p_2$ first order to be consistent with the surface-of-revolution case.

% There should be "zonal" and "highest weight" eigenfunctions on $M$. 
% \begin{itemize}
%     \item Zonal: $E_1=1,\,E_2=0$. It means $u(\cdot)$ is an eigenfunction on $M$ and $u(t,\cdot)$ is harmonic in $N$, i.e. $u(t,\phi)$ is constant in $\phi$.
%     \item Highest weight: $E_1=1,\,E_2=1$. It means $u(\cdot)$ is an eigenfunction on $M$ and $u(t,\cdot)$ is an eigenfunction on $N$.
% \end{itemize}

%    Assume $\{t=t_0\}$ is the unique critical point of $f(t)$ and $f''(t_0)\neq 0$, $f(t_0)=1$. For any sequence of $L^2(M)$-normalized, joint eigenfunction $u_h(t,\phi)$ on $M$, there is a corresponding $L^2(N)$-normalized eigenfunction $v_h(\phi)$ on $N$, such that 
%    $$
%    \|u_h\|_{L^{\infty}(M)}=Ch^{-\frac{1}{4}}\|v_h\|_{L^{\infty}(N)},
%    $$
%    where $C$ is a constant depending on $M,N,t_0$.
   
%    Conversely, any sequence of eigenfunction $v_h(\phi)$ on $N$ can be extended to a joint eigenfunction $u_h(t,\phi)$ such that $u_h(t_0,\phi)=Ch^{-\frac{1}{4}}v_h(\phi)$.

% \subsection{Manifold with a Killing field}
% Possible generalization: manifold with a Killing vector field $X$. One should take the second Hamiltonian to be such that $Op_h(p_2)=hX$. Sanity check: $[h^2\Delta_g,hX]=0$. I believe it is true because $\mathcal{L}_Xg=0$.
















%with an application of the method of steepest descent.
%This example saturates the general $O(h^{\frac{1}4})$ bound given in Theorem \ref{main1}.

%Now consider a geodesic segment \cs\cs$\gamma_3=\{(\cos\varphi, \sin\varphi, 0); \varphi \in[-\frac{\pi}4, \frac{\pi}4] \}$ which is {\em admissible}. Then
%\begin{equation*}
%\Big| \int_{\gamma_2} u_\lambda ds \Big | = \lambda^{1/4} \Big|\int_{-\frac{\pi}4}^{\frac{\pi}4} e^{i\lambda\varphi} d\varphi\Big|=\lambda^{-3/4}.
%\end{equation*}
%This bound is stronger than the general $O(h^{\frac{1}2})$ bound given in Theorem \ref{main1}.

%Let $x=(x_1, x_2)$ be geodesic normal coordinates on $\mathbb{S}^2$ at the north pole. We consider zonal harmonics centered at the north pole which can be written as oscillatory integrals of the form $u_\lambda(x)+u_\lambda(-x)$ where,
%\begin{equation}\label{zh}
%u_\lambda(x)=(2\pi \lambda)^{1/2}\int_{\mathbb{S}^1} e^{i\lambda\left<x,\omega\right>}a(x,\omega;\lambda) d\omega,
%\end{equation}
%where, $a(x,\omega;\lambda)\sim \sum_{j=0}^\infty a_j(x,\omega) \lambda^{-j}$ and $|a_0(x,\omega)|\geq \frac{1}C>0$ with $a_0(x,\omega)=1+O(|x|)$. The $\lambda^{1/2}$-factor in front of the terms in \eqref{zh} ensures that $\|u_\lambda\|_{L^2(M)}=1$. Consider a geodesic segment 
%\begin{equation*}
%\gamma_3=\{(x_1, x_2);\, x_1=\frac{\pi}2,\, x_2\in[-\frac{\pi}4, \frac{\pi}4] \}
%\end{equation*}
%which is {\em admissible}. 
%Hence one can get
%\begin{align}
%\Big| \int_{\gamma_3} u_\lambda(x) dx \Big|= & (2\pi \lambda)^{1/2} \Big| \int_{-\frac{\pi}4}^{\frac{\pi}4} \int_0^{2\pi}e^{i\lambda(\frac{\pi}2 \cos\alpha+ x_2 \sin \alpha)}a(x,\omega;\lambda) d\alpha dx_2 \Big|.
%\end{align}
%Using stationary phase in variable $(\alpha, x_2)$ gives
%\begin{equation}
%\Big| \int_{\gamma_3} u_\lambda(x) dx \Big|=O(\lambda^{-1/2}).
%\end{equation}
%This bound is consistent with the general $O(h^{\frac{1}2})$ bound given in Theorem \ref{main1}.



%Now consider another curve $\gamma_2=\{(x,y);\, x=0, y\in[0,1]\}$ which is not {\em admissible} either since \eqref{C2} is invalid. One can easily get
%\begin{equation}
%\Big| \int_{\gamma_2} \varphi_h(x,y) dy \Big |=h.
%\end{equation}

%Next two examples are motivated by \cite{}.

%For the special case of the round sphere, where
%\begin{equation*}
%\varphi_\lambda(x)=\lambda^{1/4}(x_1+i x_2)^{\lambda} 
%\end{equation*}
%are the highest-weight spherical harmonics where $\int_M |\varphi_\lambda|^2 d\text{Vol}\sim 1$, and where $\lambda=n;\quad n=1,\,2,\,3,\,\dots$. \cite{}. In terms of spherical coordinates, the restricted eigenfunction
%\begin{equation*}
%\varphi_\lambda(t)=\lambda^{1/4}[\cos t \cos \varphi+i \cos t \sin \varphi]^\lambda.
%\end{equation*}
%Consider the smooth curve $\gamma_1=\{(t, \varphi);\, t\in[-\frac{\pi}4, \frac{\pi}4],\, \varphi=0\}$ in spherical coordinate. Here $\gamma_1$ is not {\em admissible} since \eqref{C1} and \eqref{C2} are invalid. One has
%\begin{equation*}
%\Big| \int_{\gamma_1} \varphi_\lambda(t) ds \Big | \leq \lambda^{1/4} \int_{-\frac{\pi}4}^{\frac{\pi}4} |\cos t|^{\lambda} dt=O(\lambda^{-1/4})
%\end{equation*}
%with an application of the method of steepest descent.

%Now consider another curve $\gamma_2=\{(t, \varphi); t=0,\, \varphi \in[-\frac{\pi}4, \frac{\pi}4] \}$ which is {\em admissible}. Then
%\begin{equation*}
%\Big| \int_{\gamma_2} \varphi_\lambda(t) ds \Big | = \lambda^{1/4} \Big|\int_{-\frac{\pi}4}^{\frac{\pi}4} e^{i\lambda\varphi} d\varphi\Big|=\lambda^{-3/4}.
%\end{equation*}
%This bound is stronger than the general $O(h^{\frac{1}2-\delta})$ bound given in Theorem \ref{main1}.



%\cs\cs In the following we give an example where $\gamma$ is not {\em admissible}.
%Zonal harmonics.
%Let $x=(x_1, x_2)$ be geodesic normal coordinates on a convex surface of revolution centered at the north pole. We consider zonal harmonics centered at the north pole which can be written as oscillatory integrals of the form $\varphi_\lambda(x)+\varphi_\lambda(-x)$ where,
%\begin{equation}\label{zh}
%\varphi_\lambda(x)=(2\pi \lambda)^{1/2}\int_{\mathbb{S}^1} e^{i\lambda\left<x,\omega\right>}a(x,\omega;\lambda) d\omega,
%\end{equation}
%where, $a(x,\omega;\lambda)\sim \sum_{j=0}^\infty a_j(x,\omega) \lambda^{-j}$ and $|a_0(x,\omega)|\geq \frac{1}C>0$ with $a_0(x,\omega)=1+O(|x|)$. The $\lambda^{1/2}$-factor in front of the terms in \eqref{zh} ensures that $\|\varphi_\lambda\|_{L^2(M)}=1$. Consider the curve segment
%\begin{equation*}
%\gamma_1=\{(x_1, x_2);\, x_1\in[0, \pi], x_2=0 \}
%\end{equation*}
%written in the $x_j$-coordinates. Obviously $\gamma_1$ is not {\em admissible} since \eqref{C1} and \eqref{C2} are invalid. Then we can get
%\begin{align}
%\Big| \int_0^\pi \varphi_\lambda(x_1) dx_1\Big|\leq &(2\pi \lambda)^{1/2}\int_{0}^{\lambda^{-1}}\Big| \int_{\mathbb{S}^1} e^{i\lambda\left<x,\omega\right>}a(x,\omega;\lambda) d\omega\Big| dx_1 \nonumber\\
%&+(2\pi \lambda)^{1/2} \int_{\lambda^{-1}}^{\pi}\Big|\int_{\mathbb{S}^1} e^{i\lambda\left<x,\omega\right>}a(x,\omega;\lambda) d\omega\Big| dx_1 \nonumber\\
%=&O(\lambda^{-1/2})+(2\pi \lambda)^{1/2} \int_{\lambda^{-1}}^{\pi}\Big|\int_{\mathbb{S}^1} e^{i\lambda\left<x,\omega\right>}a(x,\omega;\lambda) d\omega\Big| dx_1.
%\end{align}
%Here $\left<x,\omega\right>=x_1\cos\alpha$, $\omega=(\cos\alpha, \sin\alpha)$, in the geodesic normal coordinates. Using stationary phase in $\alpha$ gives
%\begin{equation}
%\Big|\int_{\mathbb{S}^1} e^{i\lambda\left<x,\omega\right>}a(x,\omega;\lambda) d\omega\Big|=O\big((r\lambda)^{-1/2}\big).
%\end{equation}
%In conclusion
%\begin{equation}
%\Big| \int_0^\pi \varphi_\lambda(x_1)dx_1 \Big|\leq C \int_{\lambda^{-1}}^{\pi} x_1^{-1/2} dx_1+O(\lambda^{-1/2})=O(1).
%\end{equation}

%Similarly consider another curve segment, a piece of a great circle, 
%\begin{equation*}
%\gamma_2=\{(x_1, x_2);\, x_1=\frac{\pi}2,\, x_2\in[-\frac{\pi}4, \frac{\pi}4] \}
%\end{equation*}
%which is {\em admissible}. 
%Hence one can get
%\begin{align}
%\Big| \int_{\gamma_2} \varphi_\lambda(x) dx \Big|= & (2\pi \lambda)^{1/2} \Big| \int_{-\pi/4}^{\pi/4} \int_0^{2\pi}e^{i\lambda(\frac{\pi}2 \cos\alpha+ x_2 \sin \alpha)}a(x,\omega;\lambda) d\alpha dx_2 \Big|.
%\end{align}
%Using stationary phase in variable $(\alpha, x_2)$ gives
%\begin{equation}
%\Big| \int_{\gamma_2} \varphi_\lambda(x) dx \Big|=O(\lambda^{-1/2}).
%\end{equation}
%his bound is consistent with (and slightly stronger than) the general $O(h^{\frac{1}2-\delta})$ bound given in Theorem \ref{main1}.


%\cs\cs draw a picture.




%\noindent In the rest of part, we only focus on the kernel $T_1^{+}(t,x,y,\xi)$, since $T_1^{-}(t,x,y,\xi)$ can be treated similarly. For notational simplicity, we would write $T_1 (t,x,y,\xi)$ instead of $T_1^{+}(t,x,y,\xi)$.




% \subsection{More about Surfaces of revolution}
% In general, a surface of revolution in $\R^3$ is given by parametrization
% \begin{equation}\label{SoR parametrization}
%     \begin{cases}
%         x=f(t)\cos\phi\\
%         y=f(t)\sin\phi\\
%         z=g(t),
%     \end{cases}
% \end{equation}
% for $t\in [t_-,t_+]$ and $\phi\in [0,2\pi]$. Its first fundamental form is 
% \begin{align*}
%     E&=e_t\cdot e_t=f'(t)^2+g'(t)^2\\
%     F&=e_t\cdot e_\phi=0\\
%     G&=e_\phi\cdot e_\phi=f^2(t)
% \end{align*}
% and we want to arrange for $f'(t)^2+g'(t)^2=1$, so that $ds^2=dt^2+f^2(t)d\phi^2$.


% \medskip

% Is there a highest-weight eigenfunction/quasimode? For a general surface of revolution, we would guess that 
% $$
% u_{\lambda}=f(t)^{\lambda}e^{i\lambda\phi}.
% $$
% But after application of $\Delta=\partial_t^2+f^{-2}(t)\partial_\phi^2+\frac{f'(t)}{f(t)}\partial_t$ to it,
% \begin{align}
%     \Delta u_{\lambda}&=\lambda^2f'(t)^2 f(t)^{\lambda-2}e^{i\lambda\phi}-f^{-2}(t)\lambda^2f(t)^{\lambda} e^{i\lambda \phi}+lower~ order~terms\nonumber\\
%     &=-\lambda^2\Big(\frac{g'(t)^2}{f(t)^2}\Big)u_{\lambda}++lower~ order~terms.
% \end{align}
% In the sphere case, the factor $f'(t)^2/f(t)^2=1$ because $f(t)=\sin t$, $g(t)=\cos t$. In general case this is not true, and such $u_{\lambda}$ is not a quasimode. (It's a quasimode of a weighted Laplacian, but this fact does not seem useful.)

% \cs Try this.\cs Separation of variables: Let $u_{\lambda}(t,\phi)=T(t)e^{i\lambda\phi}$. 
% \begin{equation}\label{trial}
%     u_{\lambda}(t,\phi)=\exp\{\lambda\int\sqrt{\frac{|1-f^2|}{f^2}}dt\}\,e^{i\lambda \phi}.
% \end{equation}
% One can check by direct computation that 
% \begin{align}
% \Delta u_{\lambda}&=\Big(\partial_t^2+f^{-2}(t)\partial_\phi^2\Big)u_{\lambda} +lower~ order~terms\nonumber\\
% &=\Big(-\lambda^2-\lambda\frac{f'(t)}{f(t)^2\sqrt{|1-f(t)^2|}}\Big)u_{\lambda}+lower~ order~terms.
% \end{align}
% As long as $f(t)\neq 0$, it is a good quasimode. If $f(t_0)=1$, we expect that on the “equator" $f'(t_0)=0$ at correct order s.t. the second term does not blow up. \cs Make an assumption.\cs

% Back to the round sphere. $f=\cos t$, and 
% $$
% \exp\{\lambda\int\sqrt{\frac{|1-f^2|}{f^2}}dt\}=\exp\{\lambda\int |\tan t|dt\}=\exp\{\lambda \ln|\cos t|\}=(\cos t)^{\lambda}. 
% $$
% (in this convention, $t\in [-\pi/2,\pi/2]$, so $\cos t>0$.)

% \cs On warpped products, this is similar. Replace $e^{i\lambda\phi}$ by $v_{\lambda}(y)$ an eigenfunction on $N$.\cs

% The dynamics of the geodesic flow on a biaxial ellipsoid. Consider a point $p$ on the equator. By the Clairaut's relation, any geodesic staring at $p$ oscillates around the equator.
% All closed geodesic across $p$ must either be periodic or half-periodic. The latter means the first return velocity vector is symmetric with the initial velocity vector across the equator.
\bigskip
\bibliography{reference}
\bibliographystyle{alpha}

\end{document}
